\documentclass[12pt,prb,aps]{revtex4-1}
\usepackage {amsmath}
\usepackage{amssymb}
\pdfoutput = 1 
\usepackage {graphicx}
\newcommand{\bomega}{\mbox{\boldmath$\omega$}}
\allowdisplaybreaks

\begin{document}

\title{Forced Magnetic Reconnection due to Transient Resonant Magnetic Perturbations in Flowing Tokamak Plasmas}
\author{Y.~Lin and R.~Fitzpatrick\,\footnote{rfitzp@utexas.edu}}
\affiliation{Institute for Fusion Studies,  Department of Physics,  University of Texas at Austin,  Austin TX 78712, USA}

\begin{abstract}
Asymptotic matching techniques are used to investigate the linear drift-magnetohydrodynamical response of a flowing tokamak plasma to a transient  resonant magnetic perturbation (RMP). It is demonstrated that a knowledge of the solution of the ``Taylor problem," in which the RMP is suddenly switched on, and then held steady, is
sufficient to solve for the response of a transient RMP of arbitrary time dependence. 
The reconnected magnetic flux driven by the RMP  is  calculated as a function of time via  two alternative  methods.
The first solution  method, which is exact, but time consuming,  involves the numerical 
evaluation of a contour integral in Laplace transform space. The second solution method, which is approximate, but much less time consuming, involves
closed-form analytic solutions in one of seven different response regimes. 
 The results of the two solution methods are compared. 
\end{abstract}

\maketitle

\section{Introduction}
Consider a  tokamak plasma that is intrinsically stable to tearing modes, but is subject to a non-axisymmetric magnetic perturbation. 
The perturbation  inevitably 
resonates at various rational surfaces inside the plasma, and  acts to drive magnetic reconnection at these surfaces. This
process, which is known as ``forced magnetic reconnection,''  generates magnetic island chains within the plasma.\cite{err1,err1a}  Such chains degrade the plasma energy confinement,\cite{chang} may  trigger neoclassical tearing modes (NTMs),\cite{ntm2,ntm1,ntm3} and, in extreme cases, can even lead to a major disruption.

This paper  investigates the linear drift-magnetohydrodynamical (MHD) response of an idealized cylindrical  tokamak plasma to a resonant magnetic
perturbation (RMP),  resonant at a single rational surface,  that only lasts a short period of time.  Such a perturbation might be deliberately  generated by oscillating currents
flowing in external magnetic field-coils,\cite{mast,jtext}  but is much more likely to  arise spontaneously as a consequence of a sawtooth crash,\cite{ntm2} an edge localized mode
(ELMs),\cite{elm}  a fishbone, et cetera. The study of the response of a tokamak plasma to a transient RMP is particularly relevant to the problem of NTM triggering.\cite{ntm}

In this paper, we shall demonstrate that if we can predict the response of a tokamak plasma to an RMP that is suddenly switched on, and
subsequently held steady, then we  can predict the response to a transient RMP with arbitrary time variation. The problem of determining  the
response of the plasma  to a suddenly switched-on RMP, which is known as the ``Taylor problem,''  was  previously considered in a well-known paper by
Hahm \& Kulsrud,\cite{hk} and a host of subsequent papers.\cite{wb,rf1,rfb,rfc,rf2,rfd,com,hz} The present study differs from the previous studies in that it takes
plasma flows, including diamagnetic flows, into account. Such flows can profoundly modify forced magnetic reconnection. 
The ion sound-radius, and realistic levels of perpendicular energy  transport, which were previously also neglected,  are also included in the analysis. 

\section{Fundamental Analysis}

\subsection{Asymptotic Matching}
For the sake of simplicity, our tokamak equilibrium is approximated as a periodic cylinder of circular cross-section. (See Sect.~\ref{seq}.)
Suppose that the plasma equilibrium is subject to a relatively small amplitude (compared to  the equilibrium magnetic field), time-varying RMP, with $m>0$ periods in the poloidal direction, and $n>0$ periods in the
toroidal direction. The perturbation resonates with the plasma at the ``rational'' magnetic flux-surface at which the safety-factor, $q(r)$, 
takes the rational value $m/n$.\cite{err1,ric3}  Here, $r$ is a radial cylindrical coordinate. 

As far as its response to the magnetic perturbation is concerned, a high Lundquist number
tokamak plasma can be divided into two regions.\cite{fkr,hk,wb,err1,err1a,rf1,rf2,err2,com,hz,err3}
 The so-called ``inner region'' is a radially thin layer centered on the rational surface. The so-called ``outer region'' comprises everywhere in the plasma apart from the
inner region, as well as the vacuum region surrounding the plasma. 

The response of the plasma to the RMP in the outer region is
governed by the equations of linear, marginally-stable, ideal-MHD.\cite{err1} The response in the inner region is assumed to be governed by the equations of linear,
 resistive, drift-MHD.\cite{ric3} The solutions in the inner and outer regions must be asymptotically 
matched to one another.\cite{fkr}

In Refs.~\onlinecite{err3} and \onlinecite{err2}, the resonant layer equations are solved by means of a Fourier transform in space, and a {\em Fourier transform in time,}
which is appropriate to the study of non-time-dependent RMPs. However, in the present study, the layer equations are solved by means of a Fourier transform in
space, and a {\em Laplace transform in time}, which is appropriate to the study of time-dependent RMPs. Much of the analysis of Refs.~\onlinecite{err3} and \onlinecite{err2}
can be repurposed in our time-dependent analysis. However, the analysis must also be extended to some extent. [In particular, we need to calculate the
new layer parameter $F_s(\hat{g})$.] Also, the interpretation of some of the variables in our new analysis is different than in our old analysis. (In particular,
what was previously a fixed RMP frequency becomes the Laplace transform variable conjugate to time.)

\subsection{Driven Magnetic Reconnection}
Let ${\mit\Psi}_b(t)$ be the helical magnetic flux associated with the RMP to which the plasma is
subject. Likewise, let ${\mit\Psi}_0(t)$ be the driven reconnected helical magnetic flux at the rational surface. 
Making use of the analysis of Apps.~\ref{outer} and \ref{inner}, the relationship between ${\mit\Psi}_b(t)$ and ${\mit\Psi}_0(t)$ is found to be
\begin{equation}\label{drive}
{\mit\Psi}_0(\hat{t}) = \frac{1}{2\pi\,{\rm i}}\int_C \frac{E_{sb}\,F_s(\hat{g})\,\overline{\mit\Psi}_b(\hat{g})\,{\rm e}^{\,\hat{g}\,\hat{t}}}{S^{1/3}\,\hat{\mit\Delta}_s(\hat{g})-E_{ss}}\,d\hat{g},
\end{equation}
where
\begin{equation}\label{psibx}
\overline{\mit\Psi}_b(\hat{g}) = \int_0^\infty {\mit\Psi}_b(\hat{t})\,{\rm e}^{-\hat{g}\,\hat{t}}\,d\hat{t}.
\end{equation}
Here, $C$ is the Bromwich contour, and 
$\hat{t} =t/(S^{1/3}\,\tau_H)$,
whereas $S$ and $\tau_H$ are the Lundquist number and hydromagnetic time, respectively, at the rational surface. 

The real dimensionless quantities $E_{ss}$ and $E_{sb}$
are the tearing stability index, and the coupling constant between the RMP and the rational surface, respectively.\cite{ric3,spk} It is assumed that
$E_{ss}<0$, which implies that the plasma is intrinsically tearing-stable, and that any reconnected magnetic flux at the rational surface is due to the action of the RMP. 

The complex dimensionless quantities $F_s(\hat{g})$ and $\hat{\mit\Delta}_s(\hat{g})$ are determined by the resonant layer solution in the inner region. In particular,
$F_s(\hat{g})$ is the ratio of the true reconnected magnetic flux at the center of the layer to the flux at the edge of the layer,  whereas $\hat{\mit\Delta}_s(\hat{g})$ parameterizes the net helical current flowing in the  
layer. 

Note that $F_s(\hat{g})$ and $\hat{\mit\Delta}_s(\hat{g})$ depend on the following six
dimensionless  parameters that characterize the resonant layer: $Q_E= -S^{1/3}\,n\,\omega_E\,\tau_H$, $Q_e = -S^{1/3}\,n\,\omega_{\ast\,e}\,\tau_H$, $Q_i = -S^{1/3}\,n\,\omega_{\ast\,i}\,\tau_H$,
$D=S^{1/3}\,\iota_e^{\,1/2}\,\hat{d}_\beta$, $P_\perp=\tau_R/\tau_\perp$, and $P_\varphi= \tau_R/\tau_\varphi$. Here, $\tau_R$ is the resistive diffusion timescale, $\tau_\perp$ the energy confinement timescale,
$\tau_\varphi$ the momentum confinement timescale,  $\omega_E$ the ${\bf E}\times {\bf B}$ frequency, $\omega_{\ast\,e}$ the
electron diamagnetic frequency, $\omega_{\ast\,i}$ the ion diamagnetic frequency, $d_\beta=\sqrt{\beta}\,d_i$ the ion sound-radius, $d_i$ the collisionless ion skin-depth, 
$\sqrt{\beta}=\sqrt{(5/3)\,\mu_0\,p/B_\varphi^{\,2}}$ the ratio of the sound speed
to the Alfv\'{e}n speed,  $B_\varphi$ the toroidal magnetic field-strength, and $p$ the plasma pressure.   Moreover,
$P_\perp$ and $P_\varphi$ are magnetic Prandtl numbers, 
$\iota_e = - \omega_{\ast\,e}/(\omega_{\ast\,i}-\omega_{\ast\,e})$, and $\hat{d}_\beta= d_\beta/r_s$. All quantities are evaluated at the
rational surface, whose radius, $r_s$, satisfies $q(r_s)=m/n$. (See App.~\ref{symb}.)

\subsection{Taylor Problem}
Suppose that the RMP takes the idealized form 
\begin{equation}
{\mit\Psi}_b(\hat{t}) = {\mit\Psi}_{b\,0}\left\{\begin{array}{ccc} 0&~~~~&\hat{t}<0\\[0.5ex]
1&&\hat{t}\geq 0\end{array}\right..
\end{equation}
In other words, the RMP is turned on suddenly at $\hat{t}=0$, and thereafter maintains the  constant amplitude,  ${\mit\Psi}_{b\,0}$. As has already been mentioned, the problem  of determining the
response of the plasma to such an RMP as a function of time is conventionally termed the
Taylor problem (after J.B.~Taylor, who first proposed it).\cite{hk}
 It follows from Eq.~(\ref{psibx}) that
$\overline{\mit\Psi}_b(\hat{g}) ={\mit\Psi}_{b\,0}/\hat{g}$.
Let
${\mit\Psi}_0(\hat{t}) = {\mit\Psi}_\infty\,\hat{\mit\Psi}_0(\hat{t})$,
where
\begin{equation}
{\mit\Psi}_\infty = \frac{E_{sb}\,{\mit\Psi}_{b\,0}}{(-E_{ss})}
\end{equation}
is the reconnected magnetic flux that would be driven at the rational surface in the absence of any helical current flowing in the resonant layer (i.e., if $\hat{\mit\Delta}_s=0$ and $F_s=1$). 
 It follows   that
\begin{equation}\label{e29}
\hat{\mit\Psi}_0(\hat{t}) =F_{taylor}(\hat{t}), 
\end{equation}
where
\begin{equation}\label{e30x}
F_{taylor}(\hat{t}) = \frac{1}{2\pi\,{\rm i}}\int_C \frac{F_s(\hat{g})\,{\rm e}^{\,\hat{g}\,\hat{t}}}{\hat{g}\,[{\mit\Sigma}\,\hat{\mit\Delta}_s(\hat{g})+1]}\,d\hat{g},
\end{equation}
and  ${\mit\Sigma}= S^{1/3}/(-E_{ss})$. 
Note that the so-called ``Taylor function,'' $T_{taylor}(\hat{t})$, depends on the seven dimensionless layer parameters $Q_E$, $Q_e$, $Q_i$, $D$, $P_\perp$, $P_\varphi$, and ${\mit\Sigma}$. 

\subsection{Generalization of Taylor Problem}
Suppose that we have calculated the Taylor function for a particular set of layer parameters. This allows for the
immediate solution of a host of problems in which the time dependence of the RMP is more complicated than a sudden switch-on.  

For example, let our RMP be a square pulse of finite duration (in normalized time), $\tau$. In other words,
\begin{equation}\label{e145}
{\mit\Psi}_b(\hat{t}) = {\mit\Psi}_{b\,0}\left\{\begin{array}{ccl} 0&~~~~&\hat{t}<0\\[0.5ex]
1&&0 \leq \hat{t} <\tau\\[0.5ex]
0&&\hat{t}\geq \tau\end{array}\right..
\end{equation}
The response to the pulse is obviously
\begin{equation}
\hat{\mit\Psi}_0(\hat{t}) =F_{taylor}(\hat{t})- F_{taylor}(\hat{t}-\tau).
\end{equation}

Likewise, we can approximate a general RMP waveform as a histogram of the form 
\begin{equation}
{\mit\Psi}_b(\hat{t}) = {\mit\Psi}_{b\,0}\,f_i~~~~~~\hat{t}_i<\hat{t}<\hat{t}_{i+1},
\end{equation}
for $i=0,N$, where $\hat{t}_0=0$, and $\hat{t}_i<\hat{t}_{i+1}$ for all $i$. The response to the waveform is 
\begin{align}
\hat{\mit\Psi}_0(\hat{t})&=\sum_{i=0,N}f_i\left[F_{taylor}(\hat{t}-\hat{t}_i)- F_{taylor}(\hat{t}-\hat{t}_{i+1})\right]\nonumber\\[0.5ex]
&= f_0\,F_{taylor}(\hat{t}-\hat{t}_0) + \sum_{i=1,N}(f_i-f_{i-1})\,F_{taylor}(\hat{t}-\hat{t}_i) - f_N\,F_{taylor}(\hat{t}-\hat{t}_{N+1}). 
\end{align}
Thus, it is clear that if we can calculate the Taylor function then we can determine the response to a transient RMP whose time variation is arbitrary. 

\section{Numerical Approach to Problem}\label{full}
We can calculate the Taylor function, $F_{taylor}(\hat{t})$, by one of two alternative  methods.  The first solution  method,
which is exact,  involves numerically 
evaluating the contour integral (\ref{e30x}): that is, 
\begin{equation}\label{brom}
F_{taylor}(\hat{t})= \frac{1}{2\pi\,{\rm i}}\int_{\sigma-{\rm i}\,\infty}^{\sigma+{\rm i}\,\infty}\frac{F_s(\hat{g})\,{\rm e}^{\,\hat{g}\,\hat{t}}}{\hat{g}\,[{\mit\Sigma}\,\hat{\mit\Delta}_s(\hat{g})+1]}\,d\hat{g}.
\end{equation}
Here, the layer quantities $\hat{\mit\Delta}_s(\hat{g})$ and $F_s(\hat{g})$ must be calculated numerically by means of the two-step process described in App.~\ref{appd}. 
This process is a generalization of that described in App.~D of Ref.~\onlinecite{layer} that allows the calculation of $F_s(\hat{g})$, in addition to ${\mit\Delta}_s(\hat{g})$, and
is performed in Fourier space with the aid of a Ricatti transform.\cite{park}
In all of the calculations described in this paper, 
 $\sigma = 0.1$. (We have confirmed that the results of the calculation  are independent of the value of $\sigma$, provided that it is positive.) 

\begin{figure}
\centering
\includegraphics[width=0.8\textwidth]{Figure1.pdf}
\caption{\footnotesize The layer quantities $\hat{\mit\Delta}_s(\hat{g})$ and $F_s(\hat{g})$ calculated numerically along the Bromwich contour for $Q_E=0.1$,
$Q_e=1.0$, $Q_i=1.0$, $D=2.0$, $P_\perp=5.0$, $P_\varphi=5.0$, and ${\mit\Sigma}=10$.}\label{f1}
\end{figure}

Figure~\ref{f1} shows the layer quantities $\hat{\mit\Delta}_s(\hat{g})$ and $F_s(\hat{g})$ plotted along the Bromwich contour for a set of
layer parameters that are broadly consistent with those one would expect to encounter in a real tokamak plasma.\cite{err3,spk,park,lee} Figure~\ref{f2} shows the corresponding
Taylor function. It can be seen that the function is zero at $\hat{t}=0$ (i.e., at the onset of the RMP), initially grows quadratically in time, then starts to oscillate,
and eventually settles down to a steady-state value. The fact that $|F_{taylor}|\ll 1$ implies that driven magnetic reconnection is strongly suppressed by plasma
flow at the rational surface (because, in the absence of flow, we expect $F_{taylor}\rightarrow 1$ as $\hat{t}\rightarrow\infty$). 

\begin{figure}
\centering
\includegraphics[width=0.8\textwidth]{Figure2.pdf}
\caption{\footnotesize The Taylor function calculated numerically for $Q_E=0.1$,
$Q_e=1.0$, $Q_i=1.0$, $D=2.0$, $P_\perp=5.0$,  $P_\varphi=5.0$, and ${\mit\Sigma}=10$.}\label{f2}
\end{figure}

We can use the Taylor function shown in Fig.~\ref{f2} to calculate the response of the plasma at the rational surface to a transient RMP of arbitrary shape.
For example, Fig.~\ref{f3} shows the response to a square pulse of duration 20 normalized time units. As expected, the driven reconnected flux grows
when the pulse is turned on, and then decays away after it is turned off. Figure~\ref{f4} shows a more ambitious calculation that models a half sine-wave pulse as a histogram with 10 points.  The types of pulse in these two calculations could be used to model the RMP caused by 
a sawtooth crash or an ELM. 

\begin{figure}
\centering
\includegraphics[width=0.8\textwidth]{Figure3.pdf}
\caption{\footnotesize Response to a square RMP pulse calculated numerically for $Q_E=0.1$,
$Q_e=1.0$, $Q_i=1.0$, $D=2.0$, $P_\perp=5.0$, $P_\varphi=5.0$, and ${\mit\Sigma}=10$. The black curve shows the pulse.}\label{f3}
\end{figure}

\begin{figure}
\centering
\includegraphics[width=0.8\textwidth]{Figure4.pdf}
\caption{\footnotesize Response to a half sine-wave RMP pulse calculated numerically for $Q_E=0.1$,
$Q_e=1.0$, $Q_i=1.0$, $D=2.0$, $P_\perp=5.0$, $P_\varphi=5.0$, and ${\mit\Sigma}=10$. The black curve shows the pulse.}\label{f4}
\end{figure}

It should be clear, by now, that knowledge of the Taylor function for a specific resonant layer allows the  calculation of the response of the layer to
an RMP of arbitrary time variation. The only drawback to our method of calculating the Taylor function is that it is rather time consuming. Another drawback is that the
method inevitably fails at large times, because the ${\rm e}^{\,\hat{g}\,\hat{t}}$ term in the contour integral (\ref{brom}) starts to oscillate too rapidly. 
This leads us to our second  method of calculating the Taylor function, which is only approximate, but is much less time consuming than the first, and
also does not fail at large times. 

\section{Analytic Approach to Problem}\label{ana}

\subsection{Introduction}
Our second method of calculating the Taylor function involves repurposing the analysis of  Refs.~\onlinecite{err3} and \onlinecite{err2}, as described in App.~\ref{response}. 
This analysis yields analytic expressions for the layer parameters $\hat{\mit\Delta}_s(\hat{g})$ and $F_s(\hat{g})$ in seven different resonant plasma
response regimes. These expressions allow us to obtain a closed-form expression for the Taylor function in each regime using Eq.~(\ref{e30x}). The various regimes are described in the following section.

\subsection{Linear Resonant Plasma Response Regimes}\label{regimes}
According to the analysis of App.~\ref{response}, there are seven different linear resonant plasma response regimes: the inertial ($I$), the viscous-inertial ($VI$),  
the diffusive-inertial ($DI$), the resistive-inertial ($RI$), the semi-collisional ($SC$), the viscous-resistive ($VR$), and the diffusive-resistive ($DR$). These regimes can be grouped
into two classes. The four constant-$\psi$ regimes ($RI$, $SC$, $VR$, and $DR$) are characterized by $F_s\simeq 1$. In other words, a reconnected
helical flux at the center of the resonant layer that is the same as that at the edge of the layer. The three non-constant-$\psi$ regimes ($I$, $VI$, and $DI$) 
are characterized by $|F_s|\ll 1$. In other words, a reconnected
helical flux at the center of the resonant layer that is much less than that at the edge of the layer. 

The {\em inertial}\/ ($I$) regime, in which response of the resonant layer to the RMP is dominated by plasma inertia, is such that
 \begin{align}\
\hat{\mit\Delta}_s(\hat{g}) &= -\frac{\pi}{(\hat{g}+{\rm i}\,Q_E)^{1/2}\,[\hat{g}+{\rm i}\,(Q_E+Q_i)]^{1/2}},\label{in1}\\[0.5ex]
F_s(\hat{g})&= - \frac{2}{[\hat{g}+{\rm i}\,(Q_E+Q_e)]\,(\hat{g}+{\rm i}\,Q_E)\,[\hat{g}+{\rm i}\,(Q_E+Q_i)]},\label{in2}\\[0.5ex]
F_{taylor}(\hat{t}) &\simeq\frac{\epsilon_I}{{\rm i}\,(Q_E+Q_e)}
\int_0^{\hat{t}}\left[1-{\rm e}^{\,{\rm i}\,(Q_E+Q_e)\,(\hat{t}'-\hat{t})}\right]
{\rm e}^{-{\rm i}\,(Q_E+Q_i/2)\,\hat{t}'}\,J_0\!\left(\frac{Q_i}{2}\,\hat{t}'\right)d\hat{t}',\label{tinertial}\\[0.5ex]
\epsilon_I &= \frac{2}{\pi\,{\mit\Sigma}},\\[0.5ex]
\hat{\delta}(\hat{g}) &= |\hat{g}+{\rm i}\,Q_E|^{1/2}\,|\hat{g}+ {\rm i}\,(Q_E+Q_i)|^{1/2},\label{in3}\\[0.5ex]
\hat{\delta}(\hat{t})&= \frac{1}{\hat{t}}\left(1+Q_E^{\,2}\,\hat{t}^{\,2}\right)^{1/4}\left[1+(Q_E+Q_i)^2\,\hat{t}^{\,2}\right]^{1/4}.
\end{align}
Here, the radial width of the resonant layer is $\delta = r_s\,S^{-1/3}\,\hat{\delta}$,  $J_0(z)$ is
a Bessel function, and use has been
made of  Eqs.~(4.1.20)  and (5.3.34) of Ref.~\onlinecite{ed}. [The $\hat{\delta}(\hat{g})$ values come from the rule $p_\ast\sim \hat{\delta}^{-1}\sim |G|^\nu$ that is
derived in App.~\ref{response}. Obviously, $\hat{\delta}(\hat{t})$ comes from the rule $\hat{t}\sim 1/\hat{g}$.)

The {\em viscous-inertial}\/ ($VI$) regime, in which the response of the resonant layer to the RMP is dominated by plasma inertia and
viscosity, is such that
\begin{align}
\hat{\mit\Delta}_s(\hat{g}) &=- \frac{\pi}{2}\,\frac{\Gamma(1/4)}{\Gamma(3/4)}\,\frac{1}{[\hat{g}-{\rm i}\,(Q_E+Q_e)]^{1/4}\,P_\varphi^{\,1/4}},\label{vi1}\\[0.5ex]
F_s(\hat{g}) &= - \frac{\pi^{1/2}}{2}\,\frac{\Gamma(1/4)}{\Gamma(3/4)}\,\frac{1}{[\hat{g}-{\rm i}\,(Q_E+Q_e)]^{3/2}\,P_\varphi^{\,1/2}},\label{vi2}\\[0.5ex]
F_{taylor}(\hat{t})& \simeq\frac{\epsilon_{VI}}{\Gamma(5/4)}\int_0^{\hat{t}}\hat{t}'^{\,1/4}\,{\rm e}^{-{\rm i}\,(Q_E+Q_e)\,\hat{t}'}\,d\hat{t}', \\[0.5ex]
\epsilon_{VI} &= \frac{1}{\pi^{1/2}\,P_\varphi^{\,1/4}\,{\mit\Sigma}},\\[0.5ex]
\hat{\delta}(\hat{g}) &= |\hat{g}-{\rm i}\,(Q_E+Q_e)|^{1/4}\,P_\varphi^{\,1/4},\label{vi3}\\[0.5ex]
\hat{\delta}(\hat{t}) &= \frac{1}{\hat{t}^{\,1/4}}\left[1+(Q_E+Q_e)^2\,\hat{t}^{\,2}\right]^{1/8}\,P_\varphi^{\,1/4}.
\end{align}
Here, use has been made of Eq.~(5.4.1) of Ref.~\onlinecite{ed}. 

The  {\em diffusive-inertial}\/ ($DI$) regime, in which the response of the resonant layer to the RMP is dominated by plasma inertia and
cross-flux-surface energy diffusion (combined with the ion sound-radius), is such that
\begin{align}
\hat{\mit\Delta}_s(\hat{g})&= -\frac{\pi\,D}{[\hat{g}+{\rm i}\,(Q_E+Q_e)]^{1/2}\,\iota_e^{\,1/2}\,P_\perp^{\,1/2}},\label{di1}\\[0.5ex]
F_s(\hat{g}) &= - \frac{2\,D^2}{[\hat{g}+{\rm i}\,(Q_E+Q_e)]^2\,\iota_e\,P_\perp},\label{di2}\\[0.5ex]
F_{taylor}(\hat{t}) &\simeq\frac{\epsilon_{DI}}{\Gamma(3/2)}\int_0^{\hat{t}}
\hat{t}'^{\,1/2}\,{\rm e}^{-{\rm i}\,(Q_E+Q_e)\,\hat{t}'}\,d\hat{t}', \\[0.5ex]
\epsilon_{DI}& = \frac{2}{\pi}\,\frac{D}{\iota_e^{\,1/2}\,P_\perp^{\,1/2}\,{\mit\Sigma}},\\[0.5ex]
\hat{\delta}(\hat{g}) &= |\hat{g}+{\rm i}\,(Q_E+Q_e)|^{1/2}\,\iota_e^{\,1/2}\,P_\perp^{\,1/2},\label{di3}\\[0.5ex]
\hat{\delta}(\hat{t}) &=\frac{1}{\hat{t}^{\,1/2}}\left[1+(Q_E+Q_e)^2\,\hat{t}^{\,2}\right]^{1/4}\iota_e^{\,1/2}\,P_\perp^{\,1/2}.
\end{align}
Here, use has been made of has been made of Eq.~(5.4.1.) of Ref.~\onlinecite{ed}.  

The {\em resistive-inertial}\/  (RI) regime, in which the response of the resonant layer to the RMP is dominated by plasma inertia and
 resistivity, is such that
\begin{align}
\hat{\mit\Delta}_s(\hat{g}) &= \frac{2\pi\,\Gamma(3/4)}{\Gamma(1/4)}\,[\hat{g}+{\rm i}\,(Q_E+Q_e)]^{3/4}\,(\hat{g}
+{\rm i}\,Q_E)^{1/4}\,[\hat{g}+{\rm i}\,(Q_E+Q_i)]^{1/4},\label{ri1}\\[0.5ex]
F_s(\hat{g}) &\simeq 1,\\[0.5ex]
F_{taylor}(\hat{t}) &\simeq \frac{\epsilon_{RI}\,Q_i^{\,1/4}}{\sqrt{2\pi}}\int_0^{\hat{t}}{\rm e}^{-{\rm i}\,(Q_E+Q_e)\,\hat{t}'}\nonumber\\[0.5ex]
&\phantom{=}\times\left[
\int_0^{\hat{t}'}(\hat{t}'-\hat{t}'')^{-1/4}\,\hat{t}''^{-1/4}
{\rm e}^{\,{\rm i}\,(Q_e-Q_i/2)\,\hat{t}''}\,J_{-1/4}\!\left(\frac{Q_i}{2}\,\hat{t}''\right)d\hat{t}''\right]d\hat{t}',\\[0.5ex]
\epsilon_{RI}&= \frac{\Gamma(1/4)}{2\pi\,\Gamma(3/4)\,{\mit\Sigma}},\\[0.5ex]
\hat{\delta} (\hat{g})&= \frac{|\hat{g}+{\rm i}\,Q_E|^{1/4}\,|\hat{g}+{\rm i}\,(Q_E+Q_i)|^{1/4}}{|\hat{g} +{\rm i}\,(Q_E+Q_e)|^{1/4}},\label{ri3}\\[0.5ex]
\hat{\delta}(\hat{t}) &=\frac{1}{\hat{t}^{\,1/4}}\,\frac{\left(1+ Q_E^{\,2}\,\hat{t}^{\,2}\right)^{1/8}\left[1+(Q_E+Q_i)^2\,\hat{t}^{\,2}\right]^{1/8}}
{\left[1+(Q_E+Q_e)^2\,\hat{t}^{\,2}\right]^{1/8}}.
\end{align}
Here, use has been made of Eqs.~(5.4.5) and (5.4.1) of Ref.~\onlinecite{ed}, and (9.6.3) of Ref.~\onlinecite{abram}.

\begin{figure}[t]
\centering
\includegraphics[width=1.\textwidth]{Figure5.pdf}
\caption{\footnotesize Linear resonant plasma response regimes in $R$-$P$ space for the case $D = 0.9$. The various regimes are the diffusive-resistive (DR), the semi-collisional (SC), the resistive-inertial (RI), the viscous-resistive (VR), the viscous-inertial (VI), and the inertial (I).}\label{f5}
\end{figure}

The {\em semi-collisional}\/ ($SC$) regime, in which the ion sound radius plays a significant role in the response of the resonant layer to the RMP,  is such that
\begin{align}
\hat{\mit\Delta}_s(\hat{g})&= \frac{\pi\,[\hat{g}+{\rm i}\,(Q_E+Q_e)]\,(\hat{g}+{\rm i}\,Q_E)^{1/2}}{D},\label{sc1}\\[0.5ex]
F_s(\hat{g}) &\simeq 1,\\[0.5ex]
F_{taylor}(\hat{t}) &\simeq \frac{\epsilon_{SC}}{\Gamma(1/2)\,{\rm i}\,(Q_E+Q_e)} \int
_0^{\hat{t}}\left[1-{\rm e}^{\,{\rm i}\,(Q_E+Q_e)\,(\hat{t}'-\hat{t})}\right]
\hat{t}'^{-1/2}\,{\rm e}^{-{\rm i}\,Q_E\,\hat{t'}}\,d\hat{t}',\\[0.5ex]
\epsilon_{SC}&= \frac{D}{\pi\,{\mit\Sigma}},\\[0.5ex]
\hat{\delta}(\hat{g}) &= \frac{|{\hat g}+{\rm i}\,Q_E|^{1/2}}{D},\label{sc3}\\[0.5ex]
\hat{\delta}(\hat{t})& = \frac{1}{\hat{t}^{\,1/2}}\left(1+Q_E^{\,2}\,\hat{t}^{\,2}\right)^{1/4}\,D^{-1}.
\end{align}
Here, use has been made of  Eq.~(5.4.1) of Ref.~\onlinecite{ed}.

The {\em viscous-resistive}\/ ($VR$) regime, in which the response of the resonant layer to the RMP is dominated by plasma viscosity and resistivity,  is such that 
\begin{align}
\hat{\mit\Delta}_s(\hat{g}) &= \frac{6^{2/3}\,\pi\,\Gamma(5/6)}{\Gamma(1/6)}\, [\hat{g}+{\rm i}\,(Q_E+Q_e)]\,P_\varphi^{\,1/6},\label{vr1}\\[0.5ex]
F_s(\hat{g})&\simeq 1,\\[0.5ex]
F_{taylor}(\hat{t}) &= \frac{\epsilon_{VR}}{{\rm i}\,(Q_E+Q_e)+\epsilon_{VR}}\left(1 - {\rm e}^{- [{\rm i}\,(Q_E+Q_e)+\epsilon_{VR}]\,\hat{t}}\right),\\[0.5ex]
\epsilon_{VR} &= \frac{\Gamma(1/6)}{6^{2/3}\,\pi\,\Gamma(5/6)}\,\frac{1}{P_\varphi^{\,1/6}\,{\mit\Sigma}},\\[0.5ex]
\hat{\delta}(\hat{g})&=\hat{\delta}(\hat{t}) = P_\varphi^{\,1/6}.\label{vr3}
\end{align}

Finally, the {\em diffusive-resistive}\/ regime, in which the response of the resonant layer to the RMP is dominated by plasma cross-flux-surface
energy diffusion (combined with the ion sound-radius) and resistivity, is such that
\begin{align}
\hat{\mit\Delta}_s(\hat{g})&= \frac{2\pi\,\Gamma(3/4)}{\Gamma(1/4)}\,\frac{[\hat{g}+{\rm i}\,(Q_E+Q_e)]\,\iota_e^{\,1/4}\,P_\perp^{1/4}}{D^{1/2}},\label{dr1}\\[0.5ex]
F_s(\hat{g}) &\simeq 1,\\[0.5ex]
F_{taylor}(\hat{t}) &=
\frac{\epsilon_{DR}}{{\rm i}\,(Q_E+Q_e)+\epsilon_{DR}}\left(1 - {\rm e}^{- [{\rm i}\,(Q_E+Q_e)+\epsilon_{DR}]\,\hat{t}}\right),\\[0.5ex]
\epsilon_{DR} &= \frac{\Gamma(1/4)}{2\pi\,\Gamma(3/4)}\,\frac{D^{1/2}}{\iota_e^{\,1/4}\,P_\perp^{\,1/4}\,{\mit\Sigma}},\\[0.5ex]
\hat{\delta}(\hat{g}) &=\hat{\delta}(\hat{t})= \frac{\iota_e^{\,1/4}\,P_\perp^{1/4}}{D^{1/2}}.\label{dr3}
\end{align}

Note that our closed-form  expressions for the Taylor functions are all profoundly modified by plasma flow at the rational surface. In other
words, they all depend significantly on $Q_E$, $Q_e$, and $Q_i$. This dependence was missed in earlier studies.\cite{hk} However, it is
important because, in general,  we expect $|Q_E|$, $|Q_e|$, and $|Q_i|$ to all take values that are larger than, or of order, unity
in real tokamak plasmas.\cite{err3}

Suppose that $|Q_E|\sim |Q_e|\sim |Q_i|\sim Q$ and $P_\varphi\sim P_\perp\sim P$. Let $R=\sqrt{\hat{g}^2+Q^2}$. The extents of the various
linear resonant response regimes in $R$-$P$ space is illustrated in Figs.~\ref{f1} and \ref{f2}. The boundaries between the  response regimes
are taken to be the curves in parameter space at which the layer thicknesses of neighboring regimes match.
Note that Fig.~\ref{f6} differs slightly from Fig.~2 in Ref.~\onlinecite{err3} because
we are now insisting on equal layer thicknesses at the boundaries between the various response regimes. In other words, we are now of the opinion that
the boundary between the inertial and the semi-collisional regime shown in  Fig.~2 of Ref.~\onlinecite{err3}  is incorrect. 

\begin{figure}
\centering
\includegraphics[width=1.\textwidth]{Figure6.pdf}
\caption{\footnotesize Linear resonant plasma response regimes in $R$-$P$ space for the case $D = 1.2$. The various regimes are the diffusive-resistive (DR), the semi-collisional (SC), the diffusive-inertial (DI), the viscous-resistive (VR), the viscous-inertial (VI), and the inertial (I).}\label{f6}
\end{figure}

\subsection{Reconnection Scenarios}
The use of the Taylor functions for the various reconnection regimes, that are listed in Sect.~\ref{regimes}, is complicated by the fact that, as time progresses,
the resonant layer evolves from one reconnection regime to another. This evolution is illustrated in Figs.~\ref{f7} and \ref{f8}. As time passes, the
width of the resonant layer, which is initially comparatively large, gradually shrinks. Thus, on Reconnection Scenario 1 (i.e., the straight-line labelled 1
in Figs.~\ref{f7} and \ref{f8}), the resonant layer
is initially in the inertial regime. However, as soon as the inertial layer width falls below that of the viscous-inertial regime, the layer
enters the viscous-inertial regime. Likewise, as soon as the viscous-inertial layer width falls below that of the viscous-resistive regime, the
layer enters the viscous-resistive regime, where it subsequently remains. In principle, we can identify seven different reconnection scenarios (i.e., the
straight-lines labelled 1 through 7 in Figs.~\ref{f7} and \ref{f8}),
in which the layer passes through different combinations of reconnection regimes as time increases. In practice, only the scenarios characterized by $P>1$
are likely to be relevant to a real tokamak plasma.\cite{err3,spk,park,lee}

\begin{figure}[t]
\centering
\includegraphics[width=1.\textwidth]{Figure7.pdf}
\caption{Trajectory of the Fourier-Laplace transformed response of a resonant layer to an RMP in $\hat{t}$-$P$ space for the case $D=0.9$. The various regimes are the diffusive-resistive (DR), the semi-collisional (SC), the resistive-inertial (RI), the viscous-resistive (VR), the viscous-inertial (VI), and the inertial (I). }\label{f7}
\end{figure}

\begin{figure}
\centering
\includegraphics[width=1.\textwidth]{Figure8.pdf}
\caption{\footnotesize Trajectory of the Fourier-Laplace transformed response of a resonant layer to an RMP in $\hat{t}$-$P$ space for the case $D=1.2$. The various regimes are the diffusive-resistive (DR), the semi-collisional (SC), the diffusive-inertial (DI), the viscous-resistive (VR), the viscous-inertial (VI), and the inertial (I).}\label{f8}
\end{figure}

Figures~\ref{f7} and \ref{f8} indicate that the resonant layer always starts off in the inertial regime, and always ends up in either the
viscous-resistive or the diffusive-resistive regime. This suggests that if we have a very rapidly evolving, or a very transient, RMP then we can
model the resonant layer response using the inertial Taylor function, (\ref{tinertial}). On the other hand, the response to a comparatively slowly growing RMP could be modeled
using the viscous-resistive or the diffusive-resistive Taylor functions, as appropriate. We can determine exactly how rapidly or how slowly
the RMP has to evolve, in each of the two cases, by consulting Figs.~\ref{f7} and \ref{f8}.  

\subsection{Reconnection Scenario 1}
In order to perform a quick check that the approximate analytic approach to calculating the Taylor function considered in Sect.~\ref{ana}  is
compatible with the exact numerical approach  described in Sect.~\ref{full}, let us examine Reconnection Scenario 1 in more detail. The chosen layer
parameters for this scenario are $Q_E=0.1$,
$Q_e=1.0$, $Q_i=1.0$, $D=1.2$, $P_\perp=5.0$,. $P_\varphi=5.0$, and ${\mit\Sigma}=10$.

\begin{figure}[t]
\centering
\includegraphics[width=0.8\textwidth]{Figure9.pdf}
\caption{\footnotesize Taylor function at relatively early times in Reconnection Scenario 1. Here, the solid curves show the fully numerical solution, the dotted curves show the analytic solution in the inertial
response regime, and the dashed curves show the analytic solution in the viscous-inertial response regime. The two dotted vertical lines show the approximate
times at which the transitions from the inertial to the viscous-inertial, and from the viscous-inertial to the viscous-resistive, response regimes occur, in order from the left to the right. }\label{f9}
\end{figure}

Figure~\ref{f9} shows the Taylor function at relatively early times (after the RMP is switched on) in Reconnection Scenario 1. It can be seen that, in accordance with the curve labelled 1 in Fig.~\ref{f7}, the numerical function exactly matches
 the analytic function in the inertial response regime at extremely early times. At intermediate times, the numerical solution deviates from the analytic solution in the inertial regime, and approximately matches
that in the viscous-inertial response regime. 

\begin{figure}[t]
\centering
\includegraphics[width=0.8\textwidth]{Figure10.pdf}
\caption{Taylor function at relatively late times in Reconnection Scenario 1. Here, the solid curves show the numerical solution, whereas the dashed curves show the analytic solution in the 
viscous-resistive response regime. The vertical dotted line shows the approximate time at  which the transition from the viscous-inertial to the viscous-resistive response regime occurs.  }\label{f10}
\end{figure}

Figure~\ref{f10} shows the Taylor function at relatively  late times in Reconnection Scenario 1. It can be seen that,  in accordance with the curve labelled 1 in Fig.~\ref{f7}, the numerical
solution matches   the analytic solution in the viscous-resistive regime fairly well at later times.  

Figures~\ref{f9} and \ref{f10} illustrate the shortcomings  of the closed-form analytic Taylor functions. The match with the analytic viscous-inertial function is  approximate because the
system is only in the viscous-inertial response regime for a fairly limited period of time. Moreover, the match with the viscous-resistive function is not perfect because the
system is not actually in the viscous-resistive response regime at time $\hat{t}=0$, whereas the analytic function assumes that this is the case. 

The agreement between the numerical Taylor function and the analytic viscous-resistive Taylor function shown in Fig.~\ref{f10} is good enough to
prompt us to  repeat the calculation of the response of the plasma to a square RMP pulse shown in Fig.~\ref{f3} using the analytic viscous-resistive Taylor function,
rather than the numerical one. The results of the recalculation are shown in Fig.~\ref{f11}, and are surprisingly good. It can be seen that the
analytic Taylor function correctly reproduces all of  the salient features of the response, in a small fraction of the calculation time. The viscous-resistive Taylor function is the appropriate one 
because the system enters the viscous-resistive response regime after about 3.5 normalized time units, which is much less than the duration of the RMP pulse.
Incidentally, a normalized time unit would be about $5\times 10^{-5}$\,s in a real tokamak, so the pulse duration is about 1\,ms. 

\begin{figure}
\centering
\includegraphics[width=0.8\textwidth]{Figure11.pdf}
\caption{Response to a square RMP pulse calculated using the viscous-resistive analytic Taylor function for $Q_E=0.1$, $Q_e=1.0$, $P_\varphi=5.0$,
and ${\mit\Sigma}=10$. }\label{f11}
\end{figure}

\section{Summary and Discussion}
Asymptotic matching techniques are used to investigate the linear drift-MHD response of a flowing tokamak plasma to a time-varying resonant magnetic perturbation (RMP). 
It is demonstrated that a knowledge of the solution of the so-called ``Taylor problem," in which the RMP is suddenly switched on, and then held steady, is
sufficient to solve for the response of a transient RMP of arbitrary time dependence. 
The reconnected magnetic flux driven by the RMP  at the rational surface is  calculated as a function of time via  two alternative  methods.
The first solution  method, which is exact, but time consuming,  involves numerical 
evaluation of a contour integral in Laplace transform space. The second solution method, which is approximate, but much less time consuming, involves
analytic solutions in one of seven different response regimes. The extents of these response regimes in parameter space are described. 
Closed-form analytic expressions for the so-called Taylor function (i.e., the temporal solution of the Taylor problem) are obtained in each regime.
These expressions are all profoundly modified by local plasma flows.
The implementation of the second method is complicated by the fact that the system
evolves between different response regimes as time progresses. However, it is possible to represent this evolution graphically. 
 The results of the two solution methods are compared, and reasonable agreement is achieved. Moreover, in one case, the analytic Taylor function is 
 shown to be capable of reproducing all of the salient features of the response of the plasma to a transient RMP predicted by the numerical method,  at a small fraction of the calculation time. 
 
 The present study points the way to a systematic exploration of the triggering of NTMs by transient magnetic perturbations in tokamak plasmas.
 The threshold island width that a driven magnetic island chain must exceed in order to trigger an NTM is estimated to be of order the
 poloidal ion gyroradius.\cite{ntm1} Thus, we now have the tools to investigate whether the transient magnetic perturbations to which a tokamak plasma is subject are
 of sufficient amplitude and duration, and lie in the correct frequency band, to effectively seed NTMs. 

\section*{Funding}
This research was supported by the U.S.\ Department of Energy, Office of Science, Office of Fusion Energy Sciences, under  contract DE-FG02-04ER54742. 

\section*{Data Availability Statement}
The data that support the findings of this study are available from the corresponding author upon reasonable request.
\appendix

\section{Outer Region}\label{outer}
\subsection{Equilibrirum}\label{seq}
For the sake of simplicity, we shall approximate our tokamak equilibrium as a periodic cylinder of circular cross-section.\cite{ric3}
Let us employ the  conventional cylindrical coordinates, $r$, $\theta$, $z$. The
equilibrium magnetic flux-surfaces lie on surfaces of constant $r$. The system is
assumed to be periodic in the $z$ (``toroidal'') direction, with periodicity length $2\pi\,R_0$, 
where $R_0$ is the simulated  major radius of the plasma. It is helpful to define the simulated toroidal angle $\varphi=z/R_0$. Let $a$ be the minor radius of the plasma. It is assumed that $a/R_0\ll 1$. It is also assumed that the
toroidal plasma beta parameter is ${\cal O}(a^2/R_0^{\,2})$.\cite{ric3}  The equilibrium magnetic
field is written 
${\bf B} = (0,\,B_\theta(r),\, B_\varphi)$, 
  where $B_\theta(r)$ 
is the poloidal magnetic field-strength, and $B_\varphi$  the (approximately) spatially uniform  ``toroidal'' magnetic field-strength.  
The equilibrium current density takes the form
${\bf j} =(0,\,0,\,j_\varphi(r))$. The  safety-factor profile is defined as 
$q(r) = r\,B_\varphi/[R_0\,B_\theta(r)]$, and is assumed to be of order unity, and positive. 

\subsection{Cylindrical Tearing Mode Equation}
Suppose that the plasma is subject to an RMP with $m>0$ periods in the
poloidal direction, and $n>0$ periods in the toroidal direction. The perturbation is resonant at the rational surface, radius $r_s$, at which $q(r_s)=m/n$. 
This appendix considers the response of the plasma to the RMP in the outer region: that is, everywhere apart from the immediate vicinity of the rational surface. 

In the outer region, the perturbed magnetic field can be written $\delta {\bf B} = r_s\,B_\varphi\,\nabla\delta \psi\times{\bf e}_z$, where
$\delta\psi(r,\zeta,t) = \psi_1(r,t)\,{\rm e}^{\,{\rm i}\,\zeta}$,
and $\zeta=m\,\theta-n\,\varphi$ is a helical angle.\cite{err1}
The function $\psi_1(r,t)$ satisfies the ``cylindrical tearing mode equation",\cite{err1,ric3,cyl,wes}
\begin{equation}\label{e111}
\frac{\partial^2\psi_1}{\partial\hat{r}^{\,2}} + \frac{1}{\hat{r}}\,\frac{\partial\psi_1}{\partial\hat{r}}-\frac{m^2}{\hat{r}^{\,2}}\,\psi_1
- \frac{J_\varphi'\,\psi_1}{\hat{r}\,(1/q-n/m)} = 0,
\end{equation}
where $\hat{r}=r/a$, $J_\varphi'(\hat{r})= dJ_\varphi/d\hat{r}$,
$J_\varphi(\hat{r}) = R_0\,\mu_0\,j_\varphi/B_\varphi$. 
Let us assume that $J_\varphi(\hat{r}>1)=0$.  (In other words, the region $\hat{r}>1$ is assumed to be a vacuum.)
Note that Eq.~(\ref{e111}) is singular at the  rational surface. The singularity is indicative of the breakdown of the equations of linearized, 
marginally-stable, ideal-MHD in the  inner region.
 
Suppose that the plasma is surrounded by a rigid conducting wall located at $r=b$, where $b>a$. 
Let $\psi_1(b,t) = {\mit\Psi}_b(t)$, where ${\mit\Psi}_b(t)$ is a complex dimensionless quantity that specifies the amplitude and phase of the
RMP to which the plasma is subject. (Here, we are assuming that the RMP is generated by a wall perturbation, for the sake of definiteness.) The most general solution to the cylindrical tearing mode equation in the outer region is written 
\begin{equation}\label{e73}
\psi_1(\hat{r},t) = {\mit\Psi}_s(t)\,\hat\psi_s(\hat{r}) + {\mit\Psi}_b(t)\,\hat{\psi}_b(\hat{r}). 
\end{equation}
Here, $\hat{\psi}_s(\hat{r})$ is a real continuous solution of Eq.~(\ref{e111}) that satisfies 
$\hat{\psi}_s(0) = 0$, $\hat{\psi}_s(\hat{r}_s)= 1$,
and $\hat{\psi}_s(\hat{r}\geq \hat{b}) = 0$,
where $\hat{r}_s=r_s/a$, and $\hat{b}= b/a$. 
Moreover, $\hat{\psi}_b(\hat{r})$ is a real continuous solution of Eq.~(\ref{e111}) that satisfies
$\hat{\psi}_b(\hat{r}\leq \hat{r}_s) = 0$ and 
$\hat{\psi}_b(\hat{b})= 1$.\
Note that the dimensionless complex quantity ${\mit\Psi}_s(t)$ parameterizes the amplitude and phase  of the
helical magnetic flux driven  by the RMP at the interface between the inner and the outer regions.

In the limit $0<\hat{r}\ll 1$, the well-behaved solution of Eq.~(\ref{e111}) is 
$\hat{\psi}_s(\hat{r})\propto \hat{r}^{\,m}$,
Let $\hat{x} = (\hat{r}-\hat{r}_s)/\hat{r}_s$. 
In the immediate vicinity of the rational surface, the continuous solution to Eq.~(\ref{e111}), subject to the constraint 
$\hat{\psi}_s(\hat{r}_s)=1$, is 
\begin{equation}\label{e18}
\hat{\psi}_s(\hat{x}) = 1+{\mit\Delta}_{s\pm}\,\hat{x} +\alpha_s\,\hat{x}\,\ln|\hat{x}| + {\cal O}(\hat{x}^2)
\end{equation}
where $\mp$ corresponds to $\hat{x}\lessgtr 0$, $\alpha_s = -\left(\hat{r}\,q\,J_\varphi'/s\right)_{\hat{r}=\hat{r}_s}$, 
and $s=d\ln q/d\ln r$ is the magnetic shear. 
Finally, solving Eq.~(\ref{e111}) outside the plasma (where $J_\varphi'=0$), subject to the boundary condition $\hat{\psi}_s(\hat{b})=0$, yields 
\begin{equation}\label{e87}
\left(\frac{\hat{r}}{\hat{\psi}_s}\,\frac{d\hat{\psi}_s}{d\hat{r}}\right)_{\hat{r}=1} = -m\left(\frac{\hat{b}^{\,2m}+1}{\hat{b}^{\,2m}-1}\right).
\end{equation}

We can calculate $\hat{\psi}_s(\hat{r})$ by, first, launching a well-behaved solution of Eq.~(\ref{e111}) from the magnetic
axis, and integrating to just inside the rational surface.
Next, we launch a second solution of  Eq.~(\ref{e111}) from  the plasma boundary, subject to the boundary condition (\ref{e87}),   and integrate to just outside the rational surface. The two solutions are then rescaled subject to the
constraint $\hat{\psi}_s(\hat{r}_s)=1$. The real dimensionless quantity
\begin{align}\label{e88}
E_{ss} &\equiv \left[\hat{r}\,\frac{d\hat{\psi}_s}{d\hat{r}}\right]^{\hat{r}=\hat{r}_{s+}}_{\hat{r}=\hat{r}_{s-}} ={\mit\Delta}_{s+} -{\mit\Delta}_{s-}
\end{align}
is the  conventional tearing stability index.\cite{fkr,spk,ric3} It is assumed that $E_{ss}<0$, which implies that the $m/n$ tearing mode is intrinsically stable. Thus, any reconnected magnetic flux that develops at the rational surface is
due solely to the action of the RMP to which the plasma is subject.

We can calculate  $\hat{\psi}_b(\hat{r})$ by launching a solution of Eq.~(\ref{e111}) from just outside the
rational surface, such that 
$\hat{\psi}_b(\hat{r}_s) = 0$, and
$d\hat{\psi}_b(\hat{r}_s)/d\hat{r} = 1/\hat{r}_s$,
and then integrating it to  the plasma boundary. 
Let
\begin{align}
\alpha_b &= \frac{1}{2m}\left[m\,\hat{b}^{\,m}\,\hat{\psi}_b(1)+ \hat{b}^{\,m}  \,\frac{d\hat{\psi}_b(1)}{d\hat{r}}\right],\\[0.5ex]
\beta_b &= \frac{1}{2m}\left[m\,\hat{b}^{-m}\,\hat{\psi}_b(1)- \hat{b}^{-m}\,\frac{d\hat{\psi}_b(1)}{d\hat{r}}\right],
\end{align}
and
$\gamma_b= \alpha_b+\beta_b$.
The boundary condition $\hat{\psi}_b(\hat{b})=1$ is satisfied by dividing our solution by $\gamma_b$. 
Let us define the real dimensionless quantity 
\begin{equation}\label{e93}
E_{sb} \equiv \left[\hat{r}\,\frac{d\hat{\psi}_b}{d\hat{r}}\right]_{\hat{r}=\hat{r}_{s+}} = \frac{1}{\gamma_b},
\end{equation}
which parameterizes the coupling between the wall and the rational surface.\cite{spk,ric3}

\subsection{Asymptotic Matching}
The complex dimensionless quantity 
\begin{align}\label{e30}
{\mit\Delta\Psi}_{s}(t) &= \left[\hat{r}\,\frac{\partial\psi_1}{\partial\hat{r}}\right]_{\hat{r}=\hat{r}_{s-}}^{\hat{r}=\hat{r}_{s+}}
\end{align}
parameterizes the amplitude and phase of the helical current sheet that is induced in the inner region. 
Asymptotic matching between the inner and outer regions yields\,\cite{spk,ric3}
\begin{align}\label{e31a}
{\mit\Delta\Psi}_{s}(t) &= E_{ss}\,{\mit\Psi}_s(t)+E_{sb}\,{\mit\Psi}_b(t),
\end{align}
where use has been made of Eqs.~(\ref{e73}), (\ref{e88}), (\ref{e93}), and (\ref{e30}). 

Finally, Fig.~2 of Ref.~\onlinecite{spk} shows the tearing eigenfunctions, $\hat{\psi}_s(\hat{r})$ and $\hat{\psi}_b(\hat{r})$, calculated for
an $m=2/n=1$ mode in an ohmic tokamak equilibrium by the method outlined in this appendix. In this particular example, $\hat{r}_s=0.77$,  $E_{ss}= 0.90$, and $E_{sb} = 2.2$.  

\section{Inner Region}\label{inner}
\subsection{Introduction}
The inner region, which lies in the immediate vicinity of the rational surface, is assumed to be governed by the  linearized, three-field, drift-MHD model described in Ref.~\onlinecite{err3}. This
model is an extension of a model introduced in Ref.~\onlinecite{err2}, and is ultimately based on the four-field model of Ref.~\onlinecite{hkm}. 

\subsection{Linearized Three-Field Model}
Our linearized three-field model takes the form:\cite{err2,err3}
\begin{align}\label{e64}
\tau_H\,\frac{\partial\tilde{\psi}}{\partial t}-{\rm i}\,n\,(\omega_E +\omega_{\ast\,e})\,\tau_H\,\tilde{\psi} &= -{\rm i} \,\hat{x}\,(\tilde{\phi}-\tilde{N})+ S^{-1}\,\hat{\nabla}^2\tilde{\psi},\\[0.5ex]
\tau_H\,\frac{\partial\hat{\nabla}^2\tilde{\phi}}{\partial t}-{\rm i}\,n\,(\omega_E+\omega_{\ast\,i})\,\tau_H\,\hat{\nabla}^2\tilde{\phi} &= -{\rm i}\,\hat{x}\,\hat{\nabla}^2\tilde{\psi}+S^{-1}\,P_\varphi\,\hat{\nabla}^4\!\left(\tilde{\phi} + \frac{\tilde{N}}{\iota}\right),\\[0.5ex]
\tau_H\,\frac{\partial\tilde{N}}{\partial t}-{\rm i}\,n\,\omega_E\,\tau_H\,\tilde{N} &={\rm i}\,n\,\omega_{\ast\,e}\,\tau_H\,\tilde{\phi} - {\rm i}\,\iota_e\,\hat{d}_\beta^{\,2}\,\hat{x}\,\hat{\nabla}^2\tilde{\psi}+ S^{-1}\,P_\perp\,\hat{\nabla}^2 \tilde{N}.\label{e66}
\end{align}
Here, $\tilde{\psi}$ is the normalized perturbed helical magnetic flux, $\tilde{\phi}$  the normalized perturbed electrostatic potential, and $\tilde{N}$ 
the normalized perturbed plasma pressure. Moreover, $\hat{x}=(r-r_s)/r_s$, $\hat{\nabla}= r_s\,\nabla$, $\tau_H=L_s/(m\,V_A)$ is
the hydromagnetic timescale, $\omega_E=-(d{\mit\Phi}/d\psi_p)_{r_s}$ the ${\bf E}\times {\bf B}$ frequency, $\omega_{\ast\,e}=[(dp_e/d\psi_p)/(e\,n_e)]_{r_s}$
the electron diamagnetic frequency, $\omega_{\ast\,i} = - [(dp_i/d\psi_p)/(e\,n_e)]_{r_s}$ the ion diamagnetic frequency, 
$\hat{d}_\beta =\sqrt{\beta}\,d_i/r_s$ the normalized ion sound-radius, $\sqrt{\beta}=[(5/3)\,\mu_0\,(p_e+p_i)/B_\varphi^{\,2}]_{r_s}^{1/2}$ the ratio of the sound speed to the
Allfv\'{e}n speed, 
$d_i= [m_i/(n_e\,e^2\,\mu_0)]^{1/2}_{r_s}$ the collisionless ion skin-depth, 
$n_e(r)$ the equilibrium electron number density, $p_e(r)$ the equilibrium electron pressure, $p_i(r)$ the equilibrium ion pressure, ${\mit\Phi}(r)$ the equilibrium 
electrostatic potential, $\psi_p(r)= B_\varphi \int_0^r dr'\,r'/q(r')$ the equilibrium poloidal magnetic flux, 
$L_s=(R_0\,q/s)_{r_s}$ the magnetic shear length, $V_A=(B_\varphi/\sqrt{\mu_0\,n_e\,m_i})_{r_s}$ the Alfv\'{e}n speed, $S=\tau_R/\tau_H$ the Lundquist number,
$\tau_R= (\mu_0\,r^2/\eta_\parallel)_{r_s}$ the resistive diffusion timescale, $\tau_\varphi= (r_s/\chi_\varphi)_{r_s}$ the toroidal momentum confinement
time, and $\tau_\perp=(r^2/\chi_\perp)_{r_s}$ the energy confinement time. Furthermore, $e$ is the magnitude of the electron charge, $m_i$ the ion mass,
$\eta_\parallel(r)$ the plasma parallel electrical resistivity, $\chi_\varphi(r)$ the toroidal momentum diffusivity, and $\chi_\perp(r)$ the perpendicular  energy diffusivity. 
Finally, $P_\varphi=\tau_R/\tau_\varphi$ and $P_\perp=\tau_R/\tau_\perp$ are magnetic Prandtl numbers, $\iota=-\omega_{\ast\,e}/\omega_{\ast\,i}$, 
and $\iota_e=\iota/(1+\iota)$. 

 Note that we have neglected parallel ion dynamics in our model.
As is explained in Refs.~\onlinecite{err2} and
\onlinecite{err3}, this is a reasonable course of action in a low-$\beta$ plasma. 

\subsection{Laplace Transformed Layer Equations}
Let us Laplace transform Eqs.~(\ref{e64})--(\ref{e66}) by operating on them with
$\int_0^\infty (\cdots)\,{\rm e}^{-g\,t}\,dt$. Assuming that all fields are zero at $t=0$, we obtain
\begin{align}\label{e67}
g\,\tau_H\,\bar{\psi}-{\rm i}\,n\,(\omega_E +\omega_{\ast\,e})\,\tau_H\,\bar{\psi} &= -{\rm i} \,\hat{x}\,(\bar{\phi}-\bar{N})+ S^{-1}\,\hat{\nabla}^2\bar{\psi},\\[0.5ex]
g\,\tau_H\,\hat{\nabla}^2\bar{\phi}-{\rm i}\,n\,(\omega_E+\omega_{\ast\,i})\,\tau_H\,\hat{\nabla}^2\,\bar{\phi} &= -{\rm i}\,\hat{x}\,\hat{\nabla}^2\bar{\psi}+S^{-1}\,P_\varphi\,\hat{\nabla}^4\!\left(\bar{\phi} + \frac{\bar{N}}{\iota}\right),\\[0.5ex]
g\,\tau_H\,\bar{N}-{\rm i}\,n\,\omega_E\,\tau_H\,\bar{N} &= {\rm i}\,n\,\omega_{\ast\,e}\,\tau_H\,\bar{\phi} - {\rm i}\,\iota_e\,\hat{d}_\beta^{\,2}\,\hat{x}\,\hat{\nabla}^2\bar{\psi}+ S^{-1}\,P_\perp\,\hat{\nabla}^2 \bar{N},\label{e69}
\end{align}
where
\begin{equation}
\bar{\psi}(\hat{x},g) = \int_0^\infty \tilde{\psi}(\hat{x},t)\,{\rm e}^{-g\,t}\,dt,
\end{equation}
et cetera. 

\subsection{Asymptotic Matching}
Let
\begin{align}
\overline{\mit\Psi}_b(g) &= \int_0^\infty {\mit\Psi}_b(t)\,{\rm e}^{-g\,t}\,dt,\label{e71}\\[0.5ex]
\overline{\mit\Psi}_s(g) &= \int_0^\infty {\mit\Psi}_s(t)\,{\rm e}^{-g\,t}\,dt,\\[0.5ex]
\overline{\mit\Delta\Psi}_s(g) &= \int_0^\infty {\mit\Delta\Psi}_s(t)\,{\rm e}^{-g\,t}\,dt \equiv \overline{\mit\Delta}_s(g)\,\overline{\mit\Psi}_s(g).\label{e73a}
\end{align}
It follows from Eq.~(\ref{e31a}) that
\begin{equation}\label{e74}
\overline{\mit\Psi}_s(g) =\frac{ E_{sb}\,\overline{\mit\Psi}_b(g)}{\overline{\mit\Delta}_s(g)-E_{ss}}.
\end{equation}

The Laplace transformed layer equations, (\ref{e67})--(\ref{e69}), possess a tearing parity solution which is such that 
$\bar{\psi}(-\hat{x},g)=\bar{\psi}(\hat{x},g)$, $\bar{\phi}(-\hat{x},g)=-\bar{\phi}(\hat{x},g)$, and $\bar{N}(-\hat{x},g)=
-\bar{N}(\hat{x},g)$.\cite{twist} When this solution is asymptotically matched to the outer solution, we find that
\begin{equation}\label{e79}
\bar{\psi}(\hat{x},g)\rightarrow \overline{\mit\Psi}_s(g)\left[1+\overline{\mit\Delta}_s(g)\,\frac{|\hat{x}|}{2}\right]
\end{equation}
at the interface between the inner and the outer regions. 

\subsection{Stretched Radial Coordinate}
Let us define the stretched radial coordinate $X = S^{1/3}\,\hat{x}$.
Assuming that $X\sim{\cal O}(1)$ in the resonant layer that constitutes the inner region (i.e., assuming that the thickness of the layer is roughly of order $S^{-1/3}\,r_s$),
and making use of the fact that $S\gg 1$ in a conventional tokamak plasma,  Eqs.~(\ref{e67})--(\ref{e69}) reduce to the following
set of linear layer equations:
\begin{align}\label{e80}
\left[\hat{g}+{\rm i}\,(Q_E+Q_{e})\right]\bar{\psi}&= - {\rm i}\,X\left(\bar{\phi}-\bar{N}\right) + \frac{\partial^{\,2}\bar{\psi}}{\partial X^2},\\[0.5ex]
\left[\hat{g}+{\rm i}\,(Q_E +Q_{i})\right]\frac{\partial^{\,2}\bar{\phi}}{\partial X^2}&= - {\rm i}\,X\,\frac{\partial^{\,2}\bar{\psi}}{\partial X^2}+ P_\varphi\,\frac{\partial^{\,4}}{\partial X^4}\!\left(\bar{\phi} + \frac{\bar{N}}{\iota}\right),\\[0.5ex]
\left(\hat{g}+{\rm i}\,Q_E\right)\bar{N} &= - {\rm i}\,Q_{e}\,\bar{\phi}  - {\rm i} \,D^{2}\,X\,\frac{\partial^{\,2}\bar{\psi}}{\partial X^{2}}
+ P_\perp\,\frac{\partial^{\,2} \bar{N}}{\partial X^{2}},\label{e82}
\end{align}
where
$\hat{g}=S^{\,1/3}\,g\,\tau_H$,
$Q_E = -S^{\,1/3}\,n\,\omega_E\,\tau_H$, $Q_{e,i} = -S^{\,1/3}\,n\,\omega_{\ast\,e,i}\,\tau_H$,
and $D = S^{\,1/3}\,\iota_e^{1/2}\,\hat{d}_\beta$.  

According to Eq.~(\ref{e79}), the asymptotic behavior of the tearing-parity solutions of Eqs.~(\ref{e80})--(\ref{e82}) is
such that
\begin{equation}\label{e47}
\bar{\psi}(X,\hat{g})\rightarrow  \bar{\mit\Psi}_s(g)\left[\frac{\hat{\mit\Delta}_s(\hat{g})}{2}\,|X| +1 + {\cal O}\left(\frac{1}{X^2}\right)\right]
\end{equation}
as $|X|\rightarrow\infty$, where
$\overline{\mit\Delta}_s(g) = S^{1/3}\,\hat{\mit\Delta}_s(\hat{g})$.
It follows from Eqs.~(\ref{e80})--(\ref{e82}) that
\begin{equation}\label{e85}
\bar{\phi}(X,\hat{g})\rightarrow {\rm i}\,(\hat{g}+{\rm i}\,Q_E)\,\overline{\mit\Psi}_s(g)\left[ \frac{\hat{\mit\Delta}_s(\hat{g})}{2}\,{\rm sgn}(X)+ \frac{1}{X}+ {\cal O}\left(\frac{1}{X^2}\right)\right]
\end{equation}
as $|X|\rightarrow\infty$. 

\subsection{Fourier Transformation}
Equations~(\ref{e80})--(\ref{e82}) are most conveniently solved in Fourier transform space.\cite{err2} 
Let
\begin{equation}\label{e86a}
\hat{\phi}(p,\hat{g}) = \int_{-\infty}^\infty \bar{\phi}(X,\hat{g})\,{\rm e}^{-{\rm i}\,p\,X}\,dX,
\end{equation}
et cetera,
and
\begin{align}\label{e95}
Y_e(p,g)\equiv \hat{\phi}(p,\hat{g})-\hat{N}(p,\hat{g})= \pi\,[\hat{g}+{ \rm i}\,(Q_E+Q_e)]\,\overline{\mit\Psi}_s(g)\,\hat{Y}_e(p,\hat{g}).
\end{align}
The Fourier transformed layer equations can be combined to give 
\begin{align}\label{e91}
\frac{\partial}{\partial p}\!\left[A(p,\hat{g})\,\frac{\partial \hat{Y}_e}{\partial p}\right] - \frac{B(p,\hat{g})}{C(p,\hat{g})}\,p^2\,\hat{Y}_e=0,
\end{align}
where
\begin{align}
A(p,\hat{g}) &= \frac{p^2}{\hat{g}+{\rm i}\,(Q_E+Q_e)+p^2},\label{e97a}\\[0.5ex]
B(p,\hat{g})&= (\hat{g}+{\rm i}\,Q_E)\,[\hat{g}+{\rm i}\,(Q_E+Q_i)] + [\hat{g}+{\rm i}\,(Q_E+Q_i)]\,(P_\varphi+P_\perp)\,p^2 + P_\varphi\,P_\perp\,p^4,\\[0.5ex]
C(p,\hat{g}) &= \hat{g}+{\rm i}\,(Q_E+Q_e)+ \left\{P_\perp
+[\hat{g}+{\rm i}\,(Q_E+Q_i)]\,D^2\right\}p^2 + \iota_e^{\,-1}\,P_\varphi\,D^2\,p^4.\label{e97c}
\end{align}
Furthermore, because
\begin{equation}
\hat{Y}_e(p,\hat{g}) = \left[\frac{\hat{g}+{\rm i}\,(Q_E+Q_e)}{\hat{g}+{\rm i}\,Q_E}\right]\hat{\phi}(p,\hat{g})
\end{equation}
as $p\rightarrow 0$, Eqs.~(\ref{e85}) and (\ref{e95})
yield the following small-$p$ boundary condition that Eq.~(\ref{e91}) must satisfy:\,\cite{ed}
\begin{equation}\label{e101}
\hat{Y}_e(p,\hat{g})\rightarrow\frac{\hat{\mit\Delta}_s(\hat{g})}{\pi\,p} + 1+ {\cal O}(p)
\end{equation}
as $p\rightarrow 0$.  Equation~(\ref{e91}) must also satisfy the  physical boundary condition 
$\hat{Y}_e(p,\hat{g})\rightarrow 0$ 
as $p\rightarrow\infty$. 

\subsection{Reconnected Helical Magnetic Flux}\label{recon}
The true reconnected helical magnetic flux at the rational surface (divided by $r_s\,B_\varphi$)  is
\begin{equation}
{\mit\Psi}_0(\hat{t}) \equiv \psi_1(0,\hat{t}) = \tilde{\psi}(0,\hat{t}).
\end{equation}
Let 
\begin{equation}
\overline{\mit\Psi}_0(\hat{g}) = \int_0^\infty {\mit\Psi}_0(\hat{t})\,{\rm e}^{-\hat{g}\,\hat{t}}\,dt.
\end{equation}
It is easily demonstrated  that\,\cite{rf1}
\begin{equation}\label{e103}
\frac{\overline{\mit\Psi}_0(\hat{g})}{\overline{\mit\Psi}_s(\hat{g})}\equiv F_s(\hat{g}) = -\int_0^\infty A(p,\hat{g})\,\frac{\partial \hat{Y}_e}{\partial p}\,dp.
\end{equation}
Thus, it follows that
\begin{equation}\label{e110}
\overline{\mit\Psi}_0(\hat{g}) =\frac{ E_{sb}\,F_s(\hat{g})\,\overline{\mit\Psi}_b(g)}{S^{1/3}\,\hat{\mit\Delta}_s(\hat{g})-E_{ss}}.
\end{equation}

\section{Numerical Solution of Layer Equation}\label{appd}

\subsection{Introduction}
The aim of this appendix is to describe how the Fourier-Laplace transformed layer equation (\ref{e91}) can be solved numerically to
determine the two functions, $\hat{\mit\Delta}_s(\hat{g})$ and $F_s(\hat{g})$, that characterize the layer response. 

\subsection{Layer Equation}
We need to solve
\begin{equation}\label{e83}
\frac{\partial}{\partial p}\!\left(A\,\frac{\partial\hat{Y}_e}{\partial p}\right) - \frac{B}{C}\,p^2\,\hat{Y}_e,
\end{equation}
where $\hat{Y}_e(p,\hat{g})$ is the Fourier-Laplace transformed electron stream-function, and $A(p,\hat{g})$, $B(p,\hat{g})$,
and $C(p,\hat{g})$ are specified in Eqs.~(\ref{e97a})--(\ref{e97c}). 
The boundary conditions are $\hat{Y}_e(p,\hat{g})\rightarrow 0$ as $p\rightarrow\infty$, and
\begin{equation}\label{esmall}
\hat{Y}_e(p,\hat{g})\rightarrow  \frac{\hat{\mit\Delta}_s(\hat{g})}{\pi\,p}+ 1+{\cal O}(p)
\end{equation}
as $p\rightarrow 0$. We also need to calculate
\begin{equation}
F_s(\hat{g}) = -\int_0^\infty A(p,\hat{g})\,\frac{\partial Y_e}{\partial p}\,dp.
\end{equation}

Let 
\begin{equation}\label{e89}
\chi(p,\hat{g}) = A\,\frac{\partial \hat{Y}_e}{\partial p}.
\end{equation}
It follows from Eqs.~(\ref{e83}) and the previous two equations that
\begin{equation}\label{e90}
A\,\frac{\partial}{\partial p}\!\left(\frac{C}{B\,p^2}\,\frac{\partial\chi}{\partial p}\right) - \chi = 0,
\end{equation}
as well as
\begin{equation}\label{fdef}
F_s(\hat{g})= -\int_0^\infty \chi(p,\hat{g})\,dp.
\end{equation}

\subsection{Small-$p$ Limit}
The Frobenius expansion of Eq.~(\ref{e83}) about $p=0$ yields 
\begin{align}
Y_e(p,\hat{g}) &= \frac{\hat{\mit\Delta}_s}{\pi\,p} + 1
\nonumber\\[0.5ex]
&\phantom{=}+\frac{\hat{\mit\Delta}_s}{\pi}\left\{
\frac{1}{2}\,(\hat{g}+{\rm i}\,Q_E)\,[\hat{g}+{\rm i}\,(Q_E+Q_i)]-\frac{1}{\hat{g}+{\rm i}\,(Q_E+Q_e)}\right\}p\nonumber\\[0.5ex]
&\phantom{=} + \frac{1}{6}\,(\hat{g}+{\rm i}\,Q_E)\,[\hat{g}+{\rm i}\,(Q_E+Q_i)]\,p^2+{\cal O}(p^{\,3}).
\end{align}
Substituting into Eq.~(\ref{e89}), we get
\begin{align}\label{csmall}
\chi(p,\hat{g}) &= -\frac{\hat{\mit\Delta}_s}{[\hat{g}+{\rm i}\,(Q_E+Q_e)]\,\pi}+\frac{\hat{\mit\Delta}_s}{2\pi}\,\frac{(\hat{g}+{\rm i}\,Q_E)\,[g+{\rm i}\,(Q_E+Q_i)]}{\hat{g}+{\rm i}\,(Q_E+Q_e)}\,p^2\nonumber\\[0.5ex]
&\phantom{=} + \frac{1}{3} \frac{(\hat{g}+{\rm i}\,Q_E)\,[g+{\rm i}\,(Q_E+Q_i)]}{\hat{g}+{\rm i}\,(Q_E+Q_e)}\,p^3+{\cal O}(p^4).
\end{align}

\subsection{Large-$p$ Limit}
In the large-$p$ limit, if we write
\begin{align}
A &=1+\frac{\alpha}{p^2},\\[0.5ex]
\frac{B}{C} &= \beta+\frac{\gamma}{p^2},
\end{align}
and
look for a solution of Eq.~(\ref{e89}) of the form
\begin{equation}\label{elarge}
\hat{Y}_e(p,\hat{g}) \propto p^x\,\exp\left(\frac{-\beta^{1/2}\,p^2}{2}\right)
\end{equation}
then we find that
\begin{equation}
x = \frac{\gamma -\beta^{1/2}\,(1-\beta^{1/2}\,\alpha)}{2\,\beta^{1/2}}.
\end{equation}
It is easily seen that
\begin{align}
\alpha &= -[\hat{g}+ {\rm i}\,(Q_E+Q_e)],\\[0.5ex]
\beta&= \frac{\iota_e\,P_\perp}{D^2},\\[0.5ex]
\gamma&= \frac{\iota_e\,P_\perp}{D^2}\left(1+ [\hat{g}+{\rm i}\,(Q_E+Q_i)]\,\frac{P_\varphi+P_\perp}{P_\varphi\,P_\perp}
-\{P_\perp+[\hat{g}+{\rm i}\,(Q_E+Q_i)\,D^2]\}\,\frac{\iota_e}{P_\varphi\,D^2}\right).
\end{align}

Finally, it can be demonstrated from Eq.~(\ref{e89}) that
\begin{equation}\label{elargea}
\chi(p,\hat{g}) \propto p^{x+1}\,\exp\left(\frac{-\beta^{1/2}\,p^2}{2}\right)
\end{equation}
at large-$p$. 

\subsection{Ricatti Transformation}
Let 
\begin{equation}
W(p)= \frac{p}{\hat{Y}_e}\,\frac{\partial\hat{Y}_e}{\partial p}.
\end{equation}
The layer equation, (\ref{e83}), transforms to give\,\cite{park}
\begin{equation}\label{e72a}
\frac{dW}{dp} =- \frac{A'}{p}\,W -\frac{W^{\,2}}{p} + \frac{B}{A\,C}\,p^3,
\end{equation}
where
\begin{equation}
A' = \frac{\hat{g}+{\rm i}\,(Q_E+Q_e)-p^2}{\hat{g}+{\rm i}\,(Q_E+Q_e)+p^2}.
\end{equation}
This equation must be solved subject to the boundary condition that
\begin{equation}\label{e63a}
W(p) = x-\beta^{1/2}\,p^2
\end{equation}
at large-$p$ [see Eq.~(\ref{elarge})], and
\begin{equation}\label{e73ax}
W(p) =-1+\frac{\pi\,p}{\hat{\mit\Delta}_s}
\end{equation}
at small-$p$. [See Eq.~(\ref{esmall}).]

Let 
\begin{equation}\label{vdef}
V(p)= \frac{p}{\chi}\,\frac{\partial\chi}{\partial p}.
\end{equation}
Equation~(\ref{e90}) transforms to give 
\begin{equation}\label{e75a}
\frac{dV}{dp} = 2\,p\,(B'-C')\,V + \frac{3\,V}{p} - \frac{V^2}{p} + \frac{B}{A\,C}\,p^3,
\end{equation}
where 
\begin{align}
B' &= \frac{ [\hat{g}+{\rm i}\,(Q_E+Q_i)]\,(P_\varphi+P_\perp)+ 2\,P_\varphi\,P_\perp\,p^2}{ (\hat{g}+{\rm i}\,Q_E)[g+{\rm i}\,(Q_E+Q_i)] + [\hat{g}+{\rm i}\,(Q_E+Q_i)]\,(P_\varphi+P_\perp)\,p^2+ P_\varphi\,P_\perp\,p^4},\\[0.5ex]
C'&=\frac{\{P_\perp + [\hat{g}+{\rm i}\,(Q_E+Q_i)\,D^2\}+ 2\,\iota_e^{-1}\,P_\varphi\,D^2\,p^2}{\hat{g} + {\rm i}\,(Q_E+Q_e) +\{P_\perp + [\hat{g}+{\rm i}\,(Q_E+Q_i)\,D^2\}\,p^2 + \iota_e^{-1}\,P_\varphi\,D^2\,p^4}.
\end{align}
This equation must be solved subject to the boundary condition that
\begin{equation}\label{e77a}
V(p) = 1+x-\beta^{1/2}\,p^2
\end{equation}
at large-$p$. [See Eq.~(\ref{elargea}).] 

\subsection{Method of Solution}
First, we launch solutions of Eqs.~(\ref{e72a}) and (\ref{e75a}) from large $p$, subject to the respective boundary conditions (\ref{e63a}) and (\ref{e77a}), and
numerically integrate them to small $p$,  saving $V(p)$ onto a grid. We then deduce the value of $\hat{\mit\Delta}_s(\hat{g})$ from Eq.~(\ref{e73ax}). 
Second, we launch the following system of equations from small $p$, 
\begin{align}
\frac{dU}{dp} &= \frac{V(p)}{p},\\[0.5ex]
\frac{dF}{dp} &=\exp[U(p)],
\end{align}
subject to the boundary conditions
\begin{align}
U(0) &=0,\\[0.5ex]
F(0) &= 0,
\end{align}
and numerically  integrate to large $p$. Here, $V(p)$ is interpolated from the grid. 
Finally,  it follows from Eqs.~(\ref{fdef}), (\ref{csmall}), and (\ref{vdef}) that 
\begin{equation}
F_s(\hat{g}) = \frac{\hat{\mit\Delta}_s(\hat{g})\,F(\infty)}{\pi\,[\hat{g}+{\rm i}\,(Q_E+Q_e)]}.
\end{equation}

\section{Linear Resonant Response Regimes}\label{response}
\subsection{Introduction}
The aim of this section is to classify all of the different linear resonant response regimes predicted by Eqs.~(\ref{e91}), (\ref{e101}),  and (\ref{e103}). The approach taken is similar to that adopted in Refs.~\onlinecite{err3} and \onlinecite{err2},
except that we need to incorporate the normalized Laplace transform parameter, $\hat{g}$, into the analysis. We also need to calculate $F(\hat{g})$, as
well as $\hat{\mit\Delta}_s(\hat{g})$. 
In the following, we shall
assume that $|Q_E|\sim |Q_e|\sim |Q_i| \sim Q$ and $P_\varphi\sim P_\perp\sim P$ for the sake of simplicity. 

\subsection{Constant-$\psi$ Limit}\label{scpl}
Suppose that there are two layers in $p$-space. In the small-$p$ layer, suppose that Eq.~(\ref{e91}) reduces to
\begin{equation}
\frac{\partial}{\partial p}\!\left[\frac{p^2}{\hat{g}+{\rm i}\,(Q_E+Q_e)+p^2}\,\frac{\partial\hat{Y}_e}{\partial p}\right]\simeq 0
\end{equation}
when $p \sim R^{1/2}$, where $R= \sqrt{\hat{g}^2+ Q^2}$. 
 Integrating directly, making use of Eq.~(\ref{e101}), we find that
\begin{equation}\label{e107}
\hat{Y}_e(p,g)\simeq \frac{\hat{\mit\Delta}_s}{\pi}\left[\frac{1}{p} - \frac{p}{\hat{g}+{\rm i}\,(Q_E+Q_e)}\right] + 1 + {\cal O}(p^2)
\end{equation}
for $p \lesssim R^{1/2}$.

In the large-$p$ limit, for $|p|\gg {\cal O}(R^{1/2})$,  Eq.~(\ref{e91}) reduces to 
\begin{equation}\label{e329}
\frac{\partial^2\hat{Y}_e}{\partial p^2}- \frac{B(p,\hat{g})}{C(p,\hat{g})}\,p^2\,\hat{Y}_e\simeq 0,
\end{equation}
where  $\hat{Y}_e(p,\hat{g})$ must be bounded as $p\rightarrow\infty$. Asymptotic matching to the small-$p$
layer solution, (\ref{e107}), yields the following small-$p$ boundary condition  that the previous equation must satisfy:
\begin{equation}\label{e330}
\hat{Y}_e(p,\hat{g})\simeq 1-\frac{\hat{\mit\Delta}_s}{\pi}\,\frac{p}{\hat{g}+{\rm i}\,(Q_E+Q_e)} + {\cal O}(p^2)
\end{equation}
as $p\rightarrow 0$. 

In the various constant-$\psi$ linear response regimes considered in Sect.~\ref{scp}, Eq.~(\ref{e329}) reduces to an
equation of the form
\begin{equation}\label{e330a}
\frac{\partial^2\hat{Y}_e}{\partial p^2}- G\,p^k\,\hat{Y}_e\simeq 0,
\end{equation}
where $k$ is real and non-negative, and $G(\hat{g})$ is a complex constant. Let $\hat{Y}_e(p,\hat{g})=\!\sqrt{p}\,Z(q)$ and $q=\!\sqrt{G}\,p^j/j$, where
$j=(k+2)/2$. The previous equation transforms into a modified Bessel equation of general order, 
\begin{equation}
q^2\,\frac{d^2Z}{dq^2} + q\,\frac{dZ}{dq} - (q^2+\nu^2)\,Z =0,
\end{equation}
where $\nu=1/(k+2)$. The solution that is bounded as $q\rightarrow \infty$ has the small-$q$
expansion\,\cite{abram}
\begin{equation}\label{e332}
K_\nu(q)= \frac{1}{{\Gamma}(1-\nu)}\left(\frac{2}{q}\right)^\nu - \frac{1}{{\Gamma}(1+\nu)}\left(\frac{q}{2}\right)^\nu + {\cal O}\left(q^{2-\nu}\right),
\end{equation}
where $\Gamma(z)$ is a gamma function.
A comparison of the previous expression with Eq.~(\ref{e330}) reveals that
\begin{equation}
\hat{\mit\Delta}_s(\hat{g}) =  \frac{\nu^{2\nu-1}\,\pi\,\Gamma(1-\nu)}{\Gamma(\nu)}\left[\hat{g}+{\rm i}\,(Q_E+Q_e)\right]G^{\,\nu}.
\end{equation}
Note, finally, from dimensional analysis applied to Eq.~(\ref{e330a}), that $p_\ast\sim |G|^{-\nu}$, where $p_\ast$ denotes the width of the large-$p$ layer in $p$-space. This width must be larger than $R^{1/2}$ (i.e., the width of the small-$p$ layer) in order for the constant-$\psi$ approximation to hold.
Finally, it is easily demonstrated from Eqs.~(\ref{e103}), (\ref{e330}), and the fact that $p_\ast \gg R^{1/2}$, that
$F_s(\hat{g})\simeq 1$
in a constant-$\psi$ layer. Thus, as the name suggests, a constant-$\psi$ layer is characterized by a helical
magnetic flux at the center of the layer, ${\mit\Psi}_0(t)$, that is equal to that at the edge of the layer, ${\mit\Psi}_s(t)$. 

\subsection{Constant-$\psi$ Response Regimes}\label{scp}
Suppose that $R\gg P\,p^2$ and $D^2\,p^2 \ll 1$. It follows that $k=2$, $j=2$, $\nu=1/4$, and
\begin{equation}
G(\hat{g}) = \frac{(\hat{g}+{\rm i}\,Q_E)\,[\hat{g}+{\rm i}\,(Q_E+Q_i)]}{\hat{g}+{\rm i}\,(Q_E+Q_e)}.
\end{equation}
Hence, we deduce that
\begin{equation}\label{ri}
\hat{\mit\Delta}_s(\hat{g}) = \frac{2\pi\,\Gamma(3/4)}{\Gamma(1/4)}\,[\hat{g}+{\rm i}\,(Q_E+Q_e)]^{3/4}\,(\hat{g}
+{\rm i}\,Q_E)^{1/4}\,[\hat{g}+{\rm i}\,(Q_E+Q_i)]^{1/4}.
\end{equation}
This response regime is known as the {\em resistive-inertial}\/  because the layer response is dominated by
plasma resistivity and ion inertia.\cite{fkr,err1a,ara}. The characteristic layer width is $p_\ast \sim R^{-1/4}$ (which follows from the rule $p_\ast\sim |G|^{-\nu}$,
combined with $|G|\sim R$ and $\nu=1/4$),
which implies that the regime is valid when $P\ll R^{3/2}$, $R\gg D^4$, and $R\ll 1$. 

Suppose that $R\ll P\,p^2$ and $D^2\,p^2\ll 1$. It follows that $k=4$, $j=3$, $\nu=1/6$, and $G(\hat{g})=P_\varphi$. 
Hence, we deduce that
\begin{equation}\label{vr}
\hat{\mit\Delta}_s(\hat{g}) = \frac{6^{2/3}\,\pi\,\Gamma(5/6)}{\Gamma(1/6)}\, [\hat{g}+{\rm i}\,(Q_E+Q_e)]\,P_\varphi^{\,1/6}.
\end{equation}
This response regime is known as the {\em viscous-resistive}\/ regime because the layer response is dominated by
  ion perpendicular viscosity and plasma resistivity.\cite{err1,bond} The characteristic layer width is $p_\ast \sim P^{-1/6}$, which implies that
the regime is valid when  $P\gg R^{3/2}$, $P\gg D^6$, and $P\ll R^{-3}$. 

Suppose that $R\gg P\,p^2$ and $D^2\,p^2\gg 1$. It follows that $k=0$, $j=1$, $\nu=1/2$, and
\begin{equation}
G(\hat{g}) = \frac{\hat{g}+ {\rm i}\,Q_E}{D^2}.
\end{equation}
Hence, we deduce that
\begin{equation}\label{sc}
\hat{\mit\Delta}_s(\hat{g})= \frac{\pi\,[\hat{g}+{\rm i}\,(Q_E+Q_e)]\,(\hat{g}+{\rm i}\,Q_E)^{1/2}}{D}.
\end{equation}
This response regime is known as the {\em semi-collisional}\/ regime (for historical reasons).\cite{drake,wal} The ion sound-radius plays a prominent role in the
plasma response in this regime. The characteristic
layer width is $p_\ast\sim D\,R^{-1/2}$, which implies that the regime is valid when $R\ll D^4$,
 $P\ll R^2/D^2$, and $R\ll D$. 
 
Suppose, finally, that $R\ll P\,p^2$ and $D^2\,p^2\gg 1$. It follows that $k=2$, $j=2$, $\nu=1/4$, and
\begin{equation}
G(\hat{g}) = \frac{\iota_e\,P_\perp}{D^2}.
\end{equation}
Hence, we deduce that 
\begin{equation}\label{dr}
\hat{\mit\Delta}_s(\hat{g})= \frac{2\pi\,\Gamma(3/4)}{\Gamma(1/4)}\,\frac{[\hat{g}+{\rm i}\,(Q_E+Q_e)]\,\iota_e^{\,1/4}\,P_\perp^{1/4}}{D^{1/2}}.
\end{equation}
This response regime is known as the {\em diffusive-resistive}\/ regime because the layer 
response is dominated by perpendicular energy diffusivity and plasma resistivity, with the ion sound-radius also playing a role.\cite{err3} The characteristic
layer width is $p_\ast \sim P^{-1/4}\,D^{1/2}$, which implies that the regime is valid when $P\gg R^2/D^2$,
$P\ll D^6$, and $P\ll D^2/R^2$. 

\subsection{Non-constant-$\psi$ Limit}\label{sncp}
Suppose that $p\ll R^{1/2}$. In this limit, Eq.~(\ref{e91}) reduces to 
\begin{equation}
\frac{\partial}{\partial p}\!\left(p^2\,\frac{\partial\hat{Y}_e}{\partial p}\right)-[\hat{g}+{\rm i}\,(Q_E+Q_e)]\,\frac{B(p,\hat{g})}{C(p,\hat{g})}\,p^2\,\hat{Y}_e=0.
\end{equation}
In the various non-constant-$\psi$ regimes considered in Section~\ref{snonc}, the previous equation takes the form
\begin{equation}\label{ed18}
\frac{\partial}{\partial p}\!\left(p^2\,\frac{\partial \hat{Y}_e}{\partial p}\right)-G\,p^{k+2}\,\hat{Y}_e = 0,
\end{equation}
where $k$ is real and non-negative, and $G(\hat{g})$ is a complex constant. Let $V(p,\hat{g})=p\,\hat{Y}_e(p,\hat{g})$. The previous equation
yields
\begin{equation}
\frac{\partial^2 V}{\partial p^2} - G\,p^k\,V = 0.
\end{equation}
This equation is identical in form to Equation~(\ref{e330a}), which we have already solved. Indeed, the solution that
is bounded as $p\rightarrow\infty$ has the small-$p$ expansion (\ref{e332}), where $q=\!\sqrt{G}\,p^j/j$, 
$j=(k+2)/2$, and $\nu=1/(k+2)$. Matching to Equation~(\ref{e101}) yields
\begin{equation}
\hat{\mit\Delta}_s(\hat{g})= -\frac{\pi\,\nu^{1-2\nu}\,\Gamma(\nu)}{\Gamma(1-\nu)}\,\,G^{-\nu}.
\end{equation}
According to dimensional analysis applied to  Eq.~(\ref{ed18}), the layer width in $p$-space again scales as $p_\ast\sim |G|^{-\nu}$. This width must be
less that $R^{1/2}$.  
Finally, it follows from Eqs.~(\ref{e97a}) and (\ref{e103}) that
\begin{equation}
F_s(\hat{g}) = \frac{2}{\hat{g}+{\rm i}\,(Q_E+Q_e)}\int_0^\infty V(p,\hat{g})\,dp
\end{equation}
in a non-constant-$\psi$ response regime. 

\subsection{Non-constant-$\psi$ Linear Resonant Response Regimes}\label{snonc}
Suppose that $R\gg P\,p^2$ and $D^2\,p^2 \ll 1$. It follows that $k=0$, $j=1$, $\nu=1/2$, 
\begin{align}
G(\hat{g}) &= (\hat{g}+{\rm i}\,Q_E)\,[\hat{g}+{\rm i}\,(Q_E+Q_i)],\\[0.5ex]
V(p,\hat{g})&= -G^{-1/2}\,\exp(-G^{1/2}\,p).
\end{align}
Hence, we deduce that
\begin{align}\label{eia}
\hat{\mit\Delta}_s(\hat{g}) &= -\frac{\pi}{(\hat{g}+{\rm i}\,Q_E)^{1/2}\,[\hat{g}+{\rm i}\,(Q_E+Q_i)]^{1/2}},\\[0.5ex]
F_s(\hat{g})&= - \frac{2}{[\hat{g}+{\rm i}\,(Q_E+Q_e)]\,(\hat{g}+{\rm i}\,Q_E)\,[\hat{g}+{\rm i}\,(Q_E+Q_i)]}.\label{eib}
\end{align}
This response regime is known as the {\em inertial}\/ regime because the layer response is dominated by
 plasma inertia.\cite{ara,err1a} The characteristic layer width is $p_\ast \sim R^{-1}$ (which follows from the rule $p_\ast\sim |G|^\nu$, combined with $|G|\sim R^2$ and $\nu=1/2$),
which implies that the regime is valid when $P\ll R^3$, $R\gg D$, and $R\gg 1$. 

Suppose that $R\ll P\,p^2$ and $D^2\,p^2 \ll 1$. It follows that $k=2$, $j=2$, $\nu=1/4$, and
\begin{align}
G(\hat{g}) &= [\hat{g}+{\rm i}\,(Q_E+Q_e)]\,P_\varphi,\\[0.5ex]
V(p,\hat{g})&= - \frac{\Gamma(1/4)}{\sqrt{\pi}\,2^{3/4}\,G^{1/4}}\,U(0,2^{1/2}\,G^{1/4}\,p),
\end{align}
where $U(a,x)=D_{-a-1/2}(x)$ is a parabolic cylinder function.\cite{abram}
Hence, we deduce that
\begin{align}\label{evia}
\hat{\mit\Delta}_s(\hat{g}) &=- \frac{\pi}{2}\,\frac{\Gamma(1/4)}{\Gamma(3/4)}\,\frac{1}{[\hat{g}-{\rm i}\,(Q_E+Q_e)]^{1/4}\,P_\varphi^{\,1/4}},\\[0.5ex]
F_s(\hat{g}) &= - \frac{\pi^{1/2}}{2}\,\frac{\Gamma(1/4)}{\Gamma(3/4)}\,\frac{1}{[\hat{g}-{\rm i}\,(Q_E+Q_e)]^{3/2}\,P_\varphi^{\,1/2}},\label{evib}
\end{align}
where use has been made of the result
\begin{equation}
\int_0^\infty U(0,x)\,dx= \frac{\pi}{2^{3/4}\,\Gamma(3/4)}.
\end{equation}
This response regime is known as the {\em viscous-inertial}\/ regime because the layer response is dominated by
ion perpendicular viscosity and ion inertia.\cite{err1a} 
 The characteristic layer width is $p_\ast \sim R^{-1/4}\,P^{-1/4}$,
which implies that the regime is valid when $P\gg R^3$, $P\gg D^4/R$, and $P\gg R^{-3}$. 

Suppose, finally, that $R\ll P\,p^2$ and $D^2\,p^2\gg 1$. It follows that $k=0$, $j=1$, $\nu=1/2$, and
\begin{align}
G(\hat{g}) &= \frac{[\hat{g}+{\rm i}\,(Q_E+Q_e)]\,\iota_e\,P_\perp}{D^2},\\[0.5ex]
V(p,\hat{g})&= -G^{-1/2}\,\exp(-G^{1/2}\,p).
\end{align}
Hence, we deduce that
\begin{align}\label{edia}
\hat{\mit\Delta}_s&= -\frac{\pi\,D}{[\hat{g}+{\rm i}\,(Q_E+Q_e)]^{1/2}\,\iota_e^{\,1/2}\,P_\perp^{\,1/2}},\\[0.5ex]
F_s(\hat{g}) &= - \frac{2\,D^2}{[\hat{g}+{\rm i}\,(Q_E+Q_e)]^2\,\iota_e\,P_\perp}.\label{edib}
\end{align}
This response regime is known as the {\em diffusive-inertial}\/ regime because the
layer response is dominated by perpendicular energy diffusivity and ion inertia, with the ion sound-radius also playing a role.\cite{err3} The characteristic layer width is $p_\ast\sim
R^{-1/2}\,P^{-1/2}\,D$, which implies that the regime is valid when $R\ll D$, $P\ll D^4/Q$, and $P\gg D^2/R^2$. 

Note that the combination $R\gg P\,p^2$ and $D^2\,p^2\gg 1$ does not give rise to a valid non-constant-$\psi$ response regime, because
it is associated with $k=2$, which implies that $\nu$ is infinite.

\section{Table of Important Symbols}\label{symb}

\begin{tabular}{lll}
$E_{ss}$ &~~~& Tearing stability index\\[0.5ex]
$E_{sb}$ && Coupling constant between RMP and rational surface\\[0.5ex]
$S$ && Lundquist number\\[0.5ex]
$Q_E$ && Normalized ${\bf E}\times {\bf B}$ frequency\\[0.5ex]
$Q_e$ && Normalized electron diamagnetic  frequency\\[0.5ex]
$Q_i$ && Normalized ion diamagnetic frequency\\[0.5ex]
$D$ && Normalized ion sound-radius\\[0.5ex]
$P_\varphi$ && Ratio of resistive evolution time to momentum confinement time\\[0.5ex]
$P_\perp$ && Ratio of resistive evolution time to energy confinement time \\[0.5ex]
${\mit\Sigma}$ && $S^{1/3}/(-E_{ss})$\\[0.5ex]
$\hat{t}$ && Time normalized to the resistive-kink time $S^{1/3}\,\tau_H$, where $\tau_H$ is the hydrodynamic time\\[0.5ex]
${\mit\Psi}_b(\hat{t})$ && Resonant magnetic perturbation\\[0.5ex]
${\mit\Psi}_0(\hat{t})$ && Driven reconnected magnetic flux\\[0.5ex]
$T_{taylor}(\hat{t})$ && Taylor function: i.e., response to steady RMP that is switched on suddenly at $\hat{t}=0$\\[0.5ex]
$\hat{g}$ && Laplace transform variable conjugate to normalized time\\[0.5ex]
${\mit\Delta}_s(\hat{g})$ && Layer parameter specifying net current flowing in resonant layer\\[0.5ex]
$F_s(\hat{g})$ && Layer parameter specifying ratio of magnetic flux at center of  layer to that at edge.
\end{tabular}

\section*{References}
\begin{thebibliography}{99}\baselineskip 5ex

\bibitem{err1} R.~Fitzpatrick, {\em Interaction of tearing modes with external structures in cylindrical geometry (plasma)}, Nucl.\ Fusion {\bf 33}, 1049 (1993).

\bibitem{err1a} R.~Fitzpatrick,  {\em Bifurcated states of a rotating tokamak plasma in the presence of a static error-field}, Phys.\ Plasmas {\bf 5}, 3325 (1998).

\bibitem{chang}  Z.~Chang and J.D.~Callen, {\em Global energy confinement degradation due to macroscopic phenomena in tokamaks}, Nucl.\ Fusion {\bf 30}, 219 (1990).

\bibitem{ntm2} R.J.~Buttery, T.C.~Hender, D.F.~Howell, R.J.~La\,Haye, O.~Sauter, D.~Testa and EFDA-JET contributors,  {\em Onset of neoclassical tearing modes on JET}, 
Nucl.\ Fusion, {\bf 43},  69  (2003).

\bibitem{ntm1} R.J.~La\,Haye, {\em Neoclassical tearing modes and their control}, Phys.\ Plasmas {\bf 13}, 055501 (2006).

\bibitem{ntm3} R.J.~La\,Haye, C.~Chrystal, E.J.~Strait, J.D.~Callen, C.C.~Hegna, E.C.~Howell, M.~Okabayashi and R.S.~Wilcox, 
{\em Disruptive neoclassical tearing mode seeding in DIII-D with implications for ITER}, 
Nucl.\ Fusion {\bf 62}, 056017 (2022). 

\bibitem{mast} A.J.~Thornton, A.~Kirk, P.~Cahyna, I.T.~Chapman, G.~Fishpool, J.R.~Harrison, Y.Q.~Liu, L.~Kripner, M.~Peterka and 
MAST team, {\em ELM mitigation via rotating resonant magnetic perturbations on MAST}, J.\ of Nucl.\ Mat.\ {\bf 463}, 723 (2015). 

\bibitem{jtext}  D.~Li, Q.~Yu, Y.~Ding, N.~Wang, F.~Hu, R.~Jia, L.~Peng, B.~Rao, Q.~Hu, H.~Jin, et al., {\em Disruption prevention using rotating resonant magnetic perturbation on J-TEXT}, Nucl.\ Fusion {\bf 60}, 056022 (2020). 

\bibitem{elm} H.~Zohm, {\em Edge localized modes (ELMs)}, Plasma Phys.\ Controlled Fusion {\bf 38}, 105 (1996).

\bibitem{ntm} R.~Fitzpatrick, R.~Maingi, J.-K.~Park, and S.~Sabbagh
{\em Theoretical investigation of the triggering of neoclassical tearing modes by transient magnetic perturbations in NSTX}, 
Phys.\ Plasmas {\bf 30}, 072505 (2023).

\bibitem{hk} T.S.~Hahm and R.M.~Kulsrud, {\em Forced magnetic reconnection}, Phys.\ Fluids {\bf 28}, 2412 (1985).

\bibitem{wb} X.~Wang and A.~Bhattacharjee, {\em Forced reconnection and current sheet formation in Taylor’s model}, Phys.\ Fluids B {\bf 4}, 1795 (1992).

\bibitem{rf1} R.~Fitzpatrick, {\em A numerical study of forced magnetic reconnection in the viscous Taylor problem}, Phys.\ Plasmas {\bf 10}, 2304 (2003).

\bibitem{rfb} R.~Fitzpatrick, A.~Bhattacharjee, Z.W.~Ma and T.~Linde,  {\em Wave driven magnetic reconnection in the Taylor problem}, Phys.\ Plasmas {\bf 10}, 4284 (2003).

\bibitem{rfc} R. Fitzpatrick,  {\em Scaling of forced magnetic reconnection in the Hall-magnetohydrodynamic Taylor problem}, Phys.\ Plasmas {\bf 11}, 937 (2004).

\bibitem{rf2} A.~Cole and R.~Fitzpatrick, {\em Forced magnetic reconnection in the inviscid Taylor problem}, Phys.\ Plasmas {\bf 11}, 3525 (2004).

\bibitem{rfd} R.~Fitzpatrick,  {\em Scaling of forced magnetic reconnection in the Hall-magnetohydrodynamical Taylor problem with arbitrary guide field}, Phys.\ Plasmas {\bf 11}, 3961 (2004).

\bibitem{com} L.~Comisso, D.~Grasso and F.L.~Waelbroeck,  {\em Extended theory of the Taylor problem in the plasmoid-unstable regime}, Phys.\ Plasmas {\bf 22}, 042109 (2015).

\bibitem{hz} W.~Huang and P.~Zhu, {\em Analytical model of plasma response to external magnetic perturbation in absence of no-slip condition}, Phys.\  Plasmas {\bf 27}, 022514 (2020).

\bibitem{ric3} R.~Fitzpatrick, {\em Tearing mode dynamics in tokamak plasmas}, (IOP Publishing, Bristol UK, 2023). 

\bibitem{fkr} H.P.~Furth,  J.~Killeen and M.N.~Rosenbluth,  {\em Finite‐resistivity instabilities of a sheet pinch}, Phys.\ Fluids {\bf 6}, 459 (1963).

\bibitem{err3} R.~Fitzpatrick, {\em Influence of anomalous perpendicular transport on linear tearing mode dynamics in tokamak plasmas}, Phys.\ Plasmas {\bf 29}, 032507 (2022).

\bibitem{err2} A.~Cole and R.~Fitzpatrick, {\em Drift-magnetohydrodynamical model of error-field penetration in tokamak plasmas}, Phys.\ Plasmas {\bf 13}, 032503 (2006).

\bibitem{spk} R.~Fitzpatrick, {\em Error-field penetration thresholds in ohmically heated ITER and SPARC plasmas}, Phys.\ Plasmas {\bf 30}, 092512 (2023).

\bibitem{park} J.-K.~Park, {\em Parametric dependencies of resonant layer responses across linear, two-fluid, drift-MHD
regimes}, Phys.\ Plasmas {\bf 29}, 072506 (2022).

\bibitem{lee} Y.~Lee, J.-K.~Park and Y.-S.~Na, {\em Effect of parallel flow on resonant layer responses in high beta plasmas}, Nucl.\ Fusion {\bf 64}, 106058 (2024).

\bibitem{ed} A.~Erd\'{e}lyi, W.~Magnus, F.~Oberhettinger and F.G.~Tricomi, {\em Tables of integral transforms, Vol.~1}, (McGraw-Hill, New York NY, 1954).

\bibitem{abram} M.~Abramowitz and I.A.~Stegun, eds., {\em Handbook of mathematical functions, with formulas, graphs, and mathematical tables},  (Dover, New York NY, 1965).

\bibitem{cyl} H.P.~Furth, P.H.~Rutherford and H.~Selberg, {\em Tearing mode in the cylindrical tokamak}, Phys.\  Fluids {\bf 16}, 1054 (1973). 

\bibitem{layer} R.~Fitzpatrick, {\em Investigation of tearing mode stability near ideal stability boundaries via asymptotic matching techniques},  Phys.\ Plasmas {\bf 32}, 062509 (2025). 

\bibitem{wes} J.A.~Wesson, {\em Hydromagnetic stability of tokamaks}, Nucl.\ Fusion {\bf 18}, 87 (1978).

\bibitem{hkm} R.D.~Hazeltine, M.~Kotschenreuther and P.G.~Morrison, {\em A four-field model for tokamak plasma dynamics}, Phys.\ Fluids {\bf 28}, 2466 (1985).

\bibitem{twist} R.~Fitzpatrick, {\em Stability of coupled tearing and twisting modes in tokamaks}, Phys.\ Plasmas {\bf 1}, 3308 (1994).

%\bibitem{grad} I.S.~Gradshteyn and I.M.~Ryzhik, {\em Table of integrals, series, and products, corrected and enlarged edition}, (Academic Press, New York NY,  1980). 

%\bibitem{haz1} R.D.~Hazeltine and J.D.~Meiss, {\em Plasma confinement}, (Dover, New York NY, 2003).

%\bibitem{fw} R.~Fitzpatrick and F.L.~Waelbroeck, {\em Two-fluid magnetic island dynamics in slab geometry. I. Isolated islands}, Phys. Plasmas {\bf 12}, 022307 (2005).

%\bibitem{hel} R.~Fitzpatrick, {\em Helical temperature perturbations associated with tearing modes in tokamak plasmas}, Phys.\ Plasmas {\bf 2}, 825 (1995).

\bibitem{ara} G.~Ara, B.~Basu, B.~Coppi, G.~Laval, M.N.~Rosenbluth and B.V.~Waddell, {\em Magnetic reconnection and $m = 1$ oscillations in current carrying plasmas}, Ann.\ Phys.\ (NY) {\bf 112}, 443 (1978).

\bibitem{bond} A.~Bondeson and M.~Persson, {\em Resistive tearing modes in the presence of equilibrium flows}, Phys.\ Fluids {\bf 29}, 2997 (1986).

\bibitem{drake} J.F.~Drake and Y.C.~Lee,  {\em Kinetic theory of tearing instabilities}, Phys.\ Fluids {\bf 20}, 1341 (1977).

\bibitem{wal} F.L.~Waelbroeck, {\em Shielding of resonant magnetic perturbations in the long mean-free path regime}, Phys.\ Plasmas {\bf 10}, 4040 (2003).

\end{thebibliography}
\end{document}



